\exercisesection{\S~3}

\begin{enumerate}
\def\labelenumi{\arabic{enumi})}
\item
  Let G be a compact Lie group of dimension \textgreater{} 0. Then every finite commutative subgroup of G is cyclic if and only if G is isomorphic to U, SU(2, C) or the normalizer of a maximal torus in SU(2, C).
\item
  Denote by K one of the fields R, C or H, and by n an integer \(\geq 1\) . Give the space \(K_{s}^{n}\) the usual hermitian form.
\end{enumerate}

\begin{enumerate}
\def\labelenumi{\alph{enumi})}
\item
  Show that \(\mathbf{U}(n,\mathrm{K})\) is a compact real Lie group.
\item
  Show that the sphere of radius 1 in \(\mathrm{K}_s^n\) is a (real) Lie homogeneous space for \(\mathbf{U}(n,\mathbf{K})\) ; the fixer of a point is isomorphic to \(\mathbf{U}(n - 1,\mathbf{K})\) .
\item
  Deduce that the groups \(\mathbf{U}(n,\mathbf{C})\) and \(\mathbf{U}(n,\mathbf{H})\) are connected, and that \(\mathbf{O}(n,\mathbf{R})\) has two connected components.
\item
  Show that the groups \(\mathbf{SO}(n,\mathbf{R})\) and \(\mathbf{SU}(n,\mathbf{C})\) are connected.
\item
  Show that the group \(\mathbf{O}(n,\mathbf{R})\) (resp. \(\mathbf{U}(n,\mathbf{C})\) ) is the semi-direct product of Z/2Z (resp. T) by \(\mathbf{SO}(n,\mathbf{R})\) (resp. \(\mathbf{SU}(n,\mathbf{C})\) ).
\end{enumerate}

\begin{enumerate}
\def\labelenumi{\arabic{enumi})}
\setcounter{enumi}{2}
\tightlist
\item
  \begin{enumerate}
  \def\labelenumii{\alph{enumii})}
  \tightlist
  \item
    Show that the Lie algebra of the real Lie group \(\mathbf{U}(n, \mathbf{H})\) is the set of matrices \(x \in \mathrm{M}_n(\mathbf{H})\) such that \(^t\bar{x} = -x\) , with the bracket \([x, y] = xy - yx\) . Denote it by \(\mathfrak{u}(n, \mathbf{H})\) .
  \end{enumerate}
\end{enumerate}

\begin{enumerate}
\def\labelenumi{\alph{enumi})}
\setcounter{enumi}{1}
\item
  Identify \(\mathbf{C}\) with the subfield \(\mathbf{R}(i)\) of \(\mathbf{H}\) , and \(\mathbf{C}^{2n}\) with \(\mathbf{H}^n\) by the isomorphism \((z_1, \ldots, z_{2n}) \mapsto (z_1 + jz_{n+1}, \ldots, z_n + jz_{2n})\) . Prove the equality \(\mathbf{U}(n, \mathbf{H}) = \mathbf{U}(2n, \mathbf{C}) \cap \mathbf{Sp}(2n, \mathbf{C})\) .
\item
  Deduce from \(b)\) that every compact simple (real) Lie algebra of type \(\mathbf{C}_n\) is isomorphic to \(\mathfrak{u}(n, \mathbf{H})\) .
\end{enumerate}

\begin{enumerate}
\def\labelenumi{\arabic{enumi})}
\setcounter{enumi}{3}
\tightlist
\item
  Let \(n\) be an integer \(\geq 1\) .
\end{enumerate}

\begin{enumerate}
\def\labelenumi{\alph{enumi})}
\item
  Show that the group \(\mathbf{SU}(n,\mathbf{C})\) is simply-connected (use Exerc. 2).
\item
  Show that the centre of \(\mathbf{SU}(n,\mathbf{C})\) consists of the matrices \(\lambda .I_n\) for \(\lambda \in \mathbf{C}\) , \(\lambda^n = 1\) .
\item
  Every almost simple compact Lie group of type \(A_{n}\) is isomorphic to the quotient of \(\mathbf{SU}(n+1,\mathbf{C})\) by the cyclic subgroup consisting of the matrices \(\zeta^k. I_{n+1} (0 \leq k < d)\) , where \(d\) divides \(n + 1\) and \(\zeta\) is a primitive \(d\) th root of unity.
\item
  Prove that the group \(\mathbf{SL}(n,\mathbf{C})\) is simply-connected (use Chap. III, §6, no. 9, Th. 6).
\end{enumerate}

\begin{enumerate}
\def\labelenumi{\arabic{enumi})}
\setcounter{enumi}{4}
\tightlist
\item
  For \(n\) an integer \(\geq 1\) , denote by \(\mathbf{Spin}(n, \mathbf{R})\) the reduced Clifford group associated to the usual quadratic form on \(\mathbf{R}^n\) (Algebra, Chap. IX, §9, no. 5).
\end{enumerate}

\begin{enumerate}
\def\labelenumi{\alph{enumi})}
\item
  Show that \(\mathbf{Spin}(n,\mathbf{R})\) is a compact Lie group and that the surjective homomorphism \(\varphi:\mathbf{Spin}(n,\mathbf{R})\to\mathbf{SO}(n,\mathbf{R})\) (loc. cit.) is analytic, with kernel \(\{+1,-1\}\) .
\item
  Show that, for \(n \geq 2\) , \(\mathbf{Spin}(n, \mathbf{R})\) is connected and simply-connected (use Exerc. 2). The group \(\pi_{1}(\mathbf{SO}(n, \mathbf{R}))\) is cyclic of order 2.
\item
  Let \(Z_{n}\) be the centre of \(\mathbf{Spin}(n,\mathbf{R})\) . Show that \(Z_{n}=\{+1,-1\}\) if n is odd, and \(Z_{n}=\{+1,-1,\varepsilon,-\varepsilon\}\) if n is even, where \(\varepsilon=e_{1}\ldots e_{n}\) is the product of the elements of the canonical basis of \(R^{n}\) ; the group \(Z_{2r}\) is isomorphic to \((\mathbf{Z}/2\mathbf{Z})^{2}\) (resp. to \(Z/4Z\) ) if r is even (resp. odd).
\item
  Prove that every almost simple connected compact Lie group of type \(B_{n}\) ( \(n \geq 2\) ) is isomorphic to \(\mathbf{Spin}(2n + 1, \mathbf{R})\) or \(\mathbf{SO}(2n + 1, \mathbf{R})\) .
\item
  If r is odd (resp. even) and \(\geq 2\) , every almost simple connected compact Lie group of type \(D_{r}\) is isomorphic to \(\mathbf{Spin}(2r,\mathbf{R})\) , \(\mathbf{SO}(2r,\mathbf{R})\) or \(\mathbf{SO}(2r,\mathbf{R})/\{\pm I_{2r}\}\) (resp. to one of the preceding groups or \(\mathbf{Spin}(2r,\mathbf{R})/\{1,\varepsilon\}\) ).
\end{enumerate}

\begin{enumerate}
\def\labelenumi{\arabic{enumi})}
\setcounter{enumi}{5}
\tightlist
\item
  \begin{enumerate}
  \def\labelenumii{\alph{enumii})}
  \tightlist
  \item
    Show that the compact Lie group \(\mathbf{U}(n,\mathbf{H})\) is connected and simply-connected (use Exerc. 2), and that its centre is \(\{\pm I_{n}\}\) .
  \end{enumerate}
\end{enumerate}

\begin{enumerate}
\def\labelenumi{\alph{enumi})}
\setcounter{enumi}{1}
\tightlist
\item
  Every almost simple connected compact Lie group of type \(C_{n}\) ( \(n \geq 3\) ) is isomorphic to \(\mathbf{U}(n, \mathbf{H})\) or \(\mathbf{U}(n, \mathbf{H}) / \{\pm I_{n}\}\) .
\end{enumerate}

\begin{enumerate}
\def\labelenumi{\arabic{enumi})}
\setcounter{enumi}{6}
\tightlist
\item
  Let A be the algebra of Cayley octonions (Algebra, Chap. III, App., no. 3), with the basis \((e_i)_{0\leq i\leq 7}\) of loc. cit. Denote by V the subspace of pure octonions generated by \(e_1,\ldots ,e_7\) and by E the subspace of V generated by \(e_1,e_2,e_3,e_5,e_6,e_7\) . Identify the subalgebra of A generated by \(e_0,e_4\) with the field \(\mathbf{C}\) of complex numbers, and denote by G the topological group of automorphisms of the unital algebra A.
\end{enumerate}

\begin{enumerate}
\def\labelenumi{\alph{enumi})}
\item
  Denote by Q the quadratic form on V induced by the Cayley norm, so that \((e_{i})_{1\leq i\leq7}\) is an orthonormal basis of V. Prove that the map \(\sigma\mapsto\sigma|V\) is an injective homomorphism from G to the group \(\mathbf{SO}(\mathbf{Q})\) , which is isomorphic to \(\mathbf{SO}(7,\mathbf{R})\) .
\item
  Show that the multiplication of A gives E the structure of a C-vector space of dimension 3, of which \(\{e_{1}, e_{2}, e_{3}\}\) is a basis. Denote by \(\Phi\) the hermitian form on E for which this basis is orthonormal. Let H be the fixer of \(e_{4}\) in G. Show that the map \(\sigma \mapsto \sigma|E\) is an isomorphism from H to the group \(\mathbf{SU}(\Phi)\) , which is isomorphic to \(\mathbf{SU}(3, \mathbf{C})\) . The map \(\sigma \mapsto \sigma(e_{4})\) induces an embedding of G/H in the sphere of V, which is isomorphic to \(S_{6}\) .
\item
  Let T be the torus in H consisting of the automorphisms \(\sigma\) such that \(\sigma(e_{0}) = e_{0}\) , \(\sigma(e_{1}) = \alpha e_{1}\) , \(\sigma(e_{2}) = \beta e_{2}\) , \(\sigma(e_{3}) = \gamma e_{3}\) , \(\sigma(e_{4}) = e_{4}\) , \(\sigma(e_{5}) = \bar{\alpha} e_{5}\) , \(\sigma(e_{6}) = \bar{\beta} e_{6}\) , \(\sigma(e_{7}) = \bar{\gamma} e_{7}\) , where \(\alpha, \beta, \gamma\) are three complex numbers of modulus 1 such that \(\alpha\beta\gamma = 1\) . Let N be the normalizer of T in G; show that N/T is of order 12 (note that each element of N must stabilize the set of the \(\pm e_{i}, i \neq 0\) ); deduce that G is of rank 2.
\item
  Show that G is connected semi-simple of type \(G_{2}\) and that G/H can be identified with \(S_{6}\) (show that \(G_{0} \neq H\) ; deduce that \(G_{0}\) is of type \(G_{2}\) , then use a dimension argument). Every compact group of type \(G_{2}\) is isomorphic to G.
\end{enumerate}

\begin{enumerate}
\def\labelenumi{\arabic{enumi})}
\setcounter{enumi}{7}
\item
  Let G be an almost simple, connected, compact Lie group of type \(A_{n}, B_{n}, C_{n}, D_{n}\) or \(G_{2}\) . Prove that \(\pi_{2}(G) = 0\) and \(\pi_{3}(G) = Z\) (use Exerc. 2 to 7 and the fact that \(\pi_{i}(\mathbf{S}_{n})\) is zero for i \textless{} n and cyclic for i = n (cf.~General Topology, Chap. XI).
\item
  Let g be a compact Lie algebra, t a Cartan subalgebra of g; let \((\mathrm{X}_{\alpha})_{\alpha\in\mathbb{R}}\) be a Chevalley system in the split reductive algebra \((\mathfrak{g}_{\mathbf{C}},\mathfrak{t}_{\mathbf{C}})\) such that \(X_{\alpha}\) and \(X_{-\alpha}\) are conjugate (with respect to g) for all \(\alpha\in R\) .
\end{enumerate}

\begin{enumerate}
\def\labelenumi{\alph{enumi})}
\item
  Let \(\mathcal{T}\) be a \(\mathbf{Z}\) -submodule of \(\mathfrak{t}_{\mathbf{C}}\) containing the \(iH_{\alpha}\) ( \(\alpha \in \mathrm{R}\) ) and such that \(\alpha(\mathcal{T}) \subset \mathbf{Z}i\) for all \(\alpha \in \mathrm{R}\) . Show that the \(\mathbf{Z}\) -submodule \(\mathcal{G}\) of \(\mathfrak{g}_{\mathbf{C}}\) generated by \(\mathcal{T}\) and the elements \(u_{\alpha}\) and \(v_{\alpha}\) ( \(\alpha \in \mathrm{R}\) ) is a \(\mathbf{Z}\) -Lie subalgebra of \(\mathfrak{g}\) .
\item
  Assume that g (resp. t) is the Lie algebra of a compact group G (resp. of a maximal torus T of G). Let \(\Gamma(\mathrm{T})\) be the kernel of the homomorphism \(\exp_{\mathrm{T}}:t\to\mathrm{T}\) . Show that the Z-module \((2\pi)^{-1}\varGamma(\mathrm{T})\) satisfies the hypotheses of a).
\item
  Let \(\langle,\rangle\) be an invariant scalar product on g; let \(\mu\) (resp. \(\tau\) ) be the Haar measure on g (resp. t) corresponding to the Lebesgue measure when g (resp. t) is identified with a space \(R^{n}\) by means of an orthonormal basis. Denote also by \(\mu\) (resp. \(\tau\) ) the measure on g/G (resp. t/T) that is the quotient of \(\mu\) (resp. \(\tau\) ) by the normalized Haar measure on G (resp. T). Prove the formula \(\mu(\mathfrak{g}/\mathcal{G})=\tau(\mathfrak{t}/\mathcal{T}).\prod_{\alpha\in\mathrm{R}_{+}}\langle iH_{\alpha},iH_{\alpha}\rangle\) .
\end{enumerate}

\exercisesection{\S~4}

\begin{enumerate}
\def\labelenumi{\arabic{enumi})}
\tightlist
\item
  Take G to be one of the groups \(\mathbf{SU}(n,\mathbf{C})\) or \(\mathbf{U}(n,\mathbf{H})\) ; identify \(g_{C}\) with \(\mathfrak{sl}(n,\mathbf{C})\) or \(\mathfrak{sp}(2n,\mathbf{C})\) , respectively, cf.~§3, no. 4 and Exerc. 3. We use the notations of Chap. VIII, §13, nos. 1 and 3, with k=C.
\end{enumerate}

\begin{enumerate}
\def\labelenumi{\alph{enumi})}
\item
  Show that the subgroup T of G consisting of the diagonal matrices with complex coefficients is a maximal torus, and that \(\mathrm{L}(\mathrm{T})_{(\mathbf{C})} = \mathfrak{h}\) .
\item
  Identify X(T) with a subgroup of \(\mathfrak{h}^*\) by means of the homomorphism \(\delta\) . Show that the linear forms \(\varepsilon_i\) and the dominant weights \(\varpi_j\) belong to X(T). If \(t = \mathrm{diag}(t_1, \ldots, t_n) \in \mathrm{T}\) , then \(\varepsilon_i(t) = t_i\) and \(\varpi_j(t) = t_1 \ldots t_j\) for \(1 \leq i, j \leq n\) .
\item
  Deduce from b) another proof of the fact that the groups \(\mathbf{SU}(n,\mathbf{C})\) and \(\mathbf{U}(n,\mathbf{H})\) are simply-connected (cf.~§3, Exerc. 4 and 6).
\end{enumerate}

\begin{enumerate}
\def\labelenumi{\arabic{enumi})}
\setcounter{enumi}{1}
\tightlist
\item
  Take \(\mathrm{G} = \mathbf{SO}(n, \mathbf{R})\) , with \(n \geq 3\) ; put \(n = 2l + 1\) if n is odd, and n = 2l if n is even. The algebra \(g_{C}\) can be identified with \(\mathfrak{o}(n, \mathbf{C})\) ; we use the notations of Chap. VIII, §13, nos. 2 and 4. Denote by \((f_i)_{1 \leq i \leq n}\) the canonical basis of \(R^n\) . Put \(e_j = \frac{1}{\sqrt{2}}(f_{2j-1} + if_{2j})\) and \(e_{-j} = \frac{1}{\sqrt{2}}(f_{2j-1} - if_{2j})\) for \(1 \leq j \leq l\) , and \(e_0 = i\sqrt{2}f_{2l+1}\) if n is odd; choose the Witt basis of \(C^n\) given by \(e_1, \ldots, e_l, e_{-l}, \ldots, e_{-1}\) if n is even (resp. \(e_1, \ldots, e_l, e_0, e_{-l}, \ldots, e_{-1}\) if n is odd).
\end{enumerate}

\begin{enumerate}
\def\labelenumi{\alph{enumi})}
\tightlist
\item
  Let \(H_{i}\) be the subspace of \(R^{n}\) generated by \(f_{2i-1}\) and \(f_{2i}\) ( \(1 \leq i \leq l\) ); show that the subgroup of G consisting of the elements g such that \(g(\mathrm{H}_{i}) \subset \mathrm{H}_{i}\) and \(\det(g|\mathrm{H}_{i}) = 1\) for \(1 \leq i \leq l\) is a maximal torus T of G, and that \(\mathrm{L}(\mathrm{T})_{(\mathbf{C})} = \mathfrak{h}\) .
\item
  Identify X(T) with a subgroup of \(h^{*}\) by means of \(\delta\) . Show that the linear forms \(\varepsilon_{i}\) belong to X(T); if \(t \in T\) and if the restriction of t to \(H_{j}\) is a rotation through an angle \(\theta_{j}\) , then \(\varepsilon_{j}(t) = e^{i\theta_{j}}\) . The weights \(\varpi_{1}, \ldots, \varpi_{l-2}\) ; \(2\varpi_{l-1}, 2\varpi_{l}, \varpi_{l-1} \pm \varpi_{l}\) belong to X(T). If n is odd, \(\varpi_{l-1}\) belongs to X(T).
\item
  Let \(\tilde{\mathbf{G}} = \mathbf{Spin}(n, \mathbf{R})\) and let \(\varphi : \tilde{G} \to G\) be the canonical covering. For \(\boldsymbol{\theta} = (\theta_{1}, \ldots, \theta_{l}) \in \mathbf{R}^{l}\) , put \(t(\boldsymbol{\theta}) = \prod_{i=1}^{l} (\cos\theta_{i} - f_{2i-1}f_{2i}\sin\theta_{i}) \in \tilde{\mathbf{G}}\) . Show that the set of the \(t(\boldsymbol{\theta})\) for \(\theta \in R^{l}\) is a maximal torus \(\tilde{T}\) of \(\tilde{G}\) , such that \(\varphi(\tilde{T}) = T\) . If X( \(\tilde{T}\) ) is identified with a subgroup of \(h^{*}\) , we have \(\varepsilon_{j}(t(\boldsymbol{\theta})) = e^{2i\theta_{j}}\) .
\item
  Show that the weights \(\varpi_{l-1}\) and \(\varpi_{l}\) belong to X( \(\tilde{T}\) ); deduce that \(\mathbf{Spin}(n,\mathbf{R})\) is simply-connected (cf.~§3, Exerc. 5).
\end{enumerate}

\begin{enumerate}
\def\labelenumi{\arabic{enumi})}
\setcounter{enumi}{2}
\tightlist
\item
  \begin{enumerate}
  \def\labelenumii{\alph{enumii})}
  \tightlist
  \item
    Show that the automorphism \(\sigma: \mathrm{A} \mapsto \overline{\mathrm{A}}\) of \(\mathbf{SU}(n, \mathbf{C})\) is not inner for \(n \geq 3\) . Every non-inner automorphism of \(\mathbf{SU}(n, \mathbf{C})\) is of the form \((\operatorname{Int} g) \circ \sigma\) , \(g \in \mathbf{SU}(n, \mathbf{C})\) .
  \end{enumerate}
\end{enumerate}

\begin{enumerate}
\def\labelenumi{\alph{enumi})}
\setcounter{enumi}{1}
\tightlist
\item
  Show that for every almost simple connected compact group G of type \(A_{n}\) ( \(n \geq 2\) ), the group \(\operatorname{Aut}(\mathrm{G})/\operatorname{Int}(\mathrm{G})\) is cyclic of order 2 (cf.~§3, Exerc. 4).
\end{enumerate}

\begin{enumerate}
\def\labelenumi{\arabic{enumi})}
\setcounter{enumi}{3}
\tightlist
\item
  Let \(n\) be an integer \(\geq 2\) .
\end{enumerate}

\begin{enumerate}
\def\labelenumi{\alph{enumi})}
\item
  Let \(g \in \mathbf{O}(2n, \mathbf{R})\) with \(\det g = -1\) ; show that the automorphism \(\operatorname{Int} g\) of \(\mathbf{O}(2n, \mathbf{R})\) induces an automorphism of \(\mathbf{SO}(2n, \mathbf{R})\) that is not inner.
\item
  For \(n \geq 2\) , the group \(\operatorname{Aut}(\mathbf{SO}(2n, \mathbf{R}))\) is equal to \(\operatorname{Int}(\mathbf{O}(2n, \mathbf{R}))\) (which is isomorphic to \(\mathbf{O}(2n, \mathbf{R}) / \{\pm I_{2n}\}\) ).
\item
  Establish an analogous result for \(\mathbf{Spin}(2n,\mathbf{R})\) and \(\mathbf{SO}(2n,\mathbf{R})/\{\pm I_{2n}\}\) . If n is even and \(\neq 2\) , every automorphism of the group \(\mathbf{Spin}(2n,\mathbf{R})/\{1,\varepsilon\}\) (§3, Exerc. 5) is inner.
\end{enumerate}

\begin{enumerate}
\def\labelenumi{\arabic{enumi})}
\setcounter{enumi}{4}
\tightlist
\item
  Let \(\mathbf{R}\) be a reduced and irreducible root system.
\end{enumerate}

\begin{enumerate}
\def\labelenumi{\alph{enumi})}
\item
  Show that the set of roots in R of greatest length is a symmetric closed subset of R, stable under W(R).
\item
  Show that every non-empty subset P of R that is closed, symmetric and stable under W(R) is equal to R or to the set of roots of greatest length in R.
\item
  Assume that \(P \neq R\) . Show that the root system P is of type \(D_{l}\) (resp. \((\mathrm{A}_{1})^{l}, \mathrm{D}_{4}, \mathrm{A}_{2}\) ) if R is of type \(B_{l}\) (resp. \(C_{l}, F_{4}, G_{2}\) ).
\end{enumerate}

\begin{enumerate}
\def\labelenumi{\arabic{enumi})}
\setcounter{enumi}{5}
\tightlist
\item
  Let \(\mathrm{H}\) be a closed subgroup of \(\mathrm{G}\) containing \(\mathrm{N}_{\mathrm{G}}(\mathrm{T})\) .
\end{enumerate}

\begin{enumerate}
\def\labelenumi{\alph{enumi})}
\item
  Show that H is equal to its own normalizer in G, and that \(H/H_{0}\) is isomorphic to \(W/W_{H_{0}}(T)\) .
\item
  Show that \(\mathrm{R}(\mathrm{H}_{0},\mathrm{T})\) is stable under W. Conversely, if K is a connected closed subgroup containing T such that \(\mathrm{R}(\mathrm{K},\mathrm{T})\) is stable under W, the normalizer of K contains that of T.
\end{enumerate}

\(c)\) Assume that \(\mathbf{G}\) is almost simple. Prove that \(\mathbf{H}\) is equal to \(\mathrm{N_G(T)}\) , or to \(\mathbf{G}\) , or that we are (up to isomorphism) in one of the following situations:

\(\alpha\) (resp. \(\alpha'\) ) \(\mathrm{G} = \mathbf{Spin}(2l + 1, \mathbf{R})\) (resp. \(\mathrm{G} = \mathbf{SO}(2l + 1, \mathbf{R})\) ) and \(\mathrm{H}_0\) is the fixer of a non-zero vector of \(\mathbf{R}^{2l + 1}\) , isomorphic to \(\mathbf{Spin}(2l, \mathbf{R})\) (resp. \(\mathbf{SO}(2l, \mathbf{R})\) );

\(\beta)\) \(\mathrm{G} = \mathrm{U}(l,\mathbf{H})\) and \(\mathrm{H}_0\) is the subgroup D consisting of the diagonal matrices; \(\beta^{\prime})\) \(\mathrm{G} = \mathrm{U}(l,\mathbf{H}) / \{\pm 1\}\) and \(\mathrm{H}_0 = \mathrm{D} / \{\pm 1\}\) ;

\(\gamma)\) G is of type \(\mathrm{F}_4\) and \(\mathrm{H}_0\) is isomorphic to Spin(8,R);

\(\delta)\) G is of type \(\mathrm{G}_2\) and \(\mathrm{H}_0\) is the subgroup (isomorphic to \(\mathbf{SU}(3,\mathbf{C})\) ) defined in Exerc. 7 of §3.

\begin{enumerate}
\def\labelenumi{\arabic{enumi})}
\setcounter{enumi}{6}
\tightlist
\item
  Let \(\tau : \mathbf{G} \to \mathbf{GL}(\mathbf{V})\) be a continuous representation of \(\mathbf{G}\) on a finite dimensional real vector space. Assume that for all \(\lambda \in \mathrm{X}(\mathrm{T})\) , we have \(\dim_{\mathbf{C}} \hat{V}_{\lambda} \leq 1\) .
\end{enumerate}

\begin{enumerate}
\def\labelenumi{\alph{enumi})}
\item
  Show that the representation \(\tau\) is the direct sum of a finite family \((\tau_{i})_{1\leq i\leq s}\) of irreducible representations, mutually non-isomorphic, whose commutants \(K_{i}\ (1\leq i\leq s)\) are isomorphic to R or C.
\item
  There exists a C-vector space structure on V for which the operations of \(\tau(\mathrm{G})\) are C-linear if and only if \(K_{1},\ldots,K_{s}\) are isomorphic to C; there are then \(2^{s}\) structures of this type.
\end{enumerate}

\begin{enumerate}
\def\labelenumi{\arabic{enumi})}
\setcounter{enumi}{7}
\tightlist
\item
  Let H be a connected closed subgroup of G of maximum rank, h its Lie algebra, X the manifold G/H and V the tangent space of X at the point corresponding to the class of H; identify V with the quotient g/h.
\end{enumerate}

\begin{enumerate}
\def\labelenumi{\alph{enumi})}
\tightlist
\item
  If j is an almost complex structure (Differentiable and Analytic Manifolds, Results, 8.8.3) on X, denote by \(V''(j)\) the subspace of the complex vector space \(C \otimes V\) consisting of the elements u such that \(j(u) = -iu\) , and by \(\mathfrak{q}(j)\) the subspace of \(g_{C}\) that is the inverse image of \(V''(j)\) , so that the canonical map \(V \to \mathfrak{g}_{\mathbf{C}} / \mathfrak{q}(j)\) is a C-linear isomorphism when V is given the C-vector space structure induced by j. Prove that the map \(j \mapsto \mathfrak{q}(j)\) is a bijection from the set of almost complex structures on X invariant under G to the set of complex subspaces p of \(g_{C}\) satisfying the following conditions:
\end{enumerate}

\[
\mathfrak {p} + \bar {\mathfrak {p}} = \mathfrak {g} _ {\mathbf {C}}\tag{1}
\]

\[
\mathfrak {p} \cap \bar {\mathfrak {p}} = \mathfrak {h} _ {\mathbf {C}}\tag{2}
\]

\[
[ \mathfrak {h}, \mathfrak {p} ] \subset \mathfrak {p}.\tag{3}
\]

\begin{enumerate}
\def\labelenumi{\alph{enumi})}
\setcounter{enumi}{1}
\item
  There exists such a structure on X if and only if the commuting fields of the irreducible subrepresentations of the adjoint representation of H on V are all isomorphic to C; in that case, there are \(2^{s}\) such structures, where s is the number of irreducible subrepresentations of V (use Exerc. 7).
\item
  Let j be an almost complex structure on X, invariant under G, and \(\mathfrak{p} = \mathfrak{q}(j)\) the associated subspace. Then, j is integrable (that is, associated to a complex-analytic manifold structure on X, cf.~Differentiable and Analytic Manifolds, Results, 8.8.5 to 8.8.8) if and only if p satisfies the condition
\end{enumerate}

\[
[ \mathfrak {p}, \mathfrak {p} ] \subset \mathfrak {p}.\tag{4}
\]

\begin{enumerate}
\def\labelenumi{\alph{enumi})}
\setcounter{enumi}{3}
\item
  There exists a complex structure (that is, an integrable almost complex structure) on X invariant under G if and only if H is the centralizer of a torus in G (show that conditions (1) to (4) above imply that p is a parabolic subalgebra of \(g_{C}\) (Chap. VIII, §3, no. 5)); in that case, these complex structures correspond bijectively to the parabolic subalgebras p of \(g_{C}\) that are the direct sum of \(h_{C}\) and their unipotent radical (cf.~Chap. VIII, §3, no. 4).
\item
  Prove that there exist exactly \(\operatorname{Card}(W)\) complex structures on G/T invariant under G; if \(\sigma\) and \(\sigma'\) are two such structures, there exists a unique element \(w \in W\) such that the canonical operation of w on G/T (by inner automorphisms) transforms \(\sigma\) into \(\sigma'\) . If w is a non-identity element of W, the operation of w on G/T is not C-analytic for any complex structure on G/T invariant under G.
\item
  Determine the complex structures on \(S^{2}\) invariant under \(\mathbf{SO}(3)\) .
\end{enumerate}

\(g)^{*}\) With the notations of Exerc. 8 \(d\) ), let \(\mathrm{G}_c\) be the complexification of \(\mathrm{G}\) and \(\mathrm{P}\) the complex Lie subgroup of \(\mathrm{G}_c\) with Lie algebra \(\mathfrak{p}\) . Show that the canonical map \(\mathrm{G / H} \to \mathrm{G}_c / \mathrm{P}\) is an isomorphism of complex-analytic manifolds.*

\begin{enumerate}
\def\labelenumi{\arabic{enumi})}
\setcounter{enumi}{8}
\tightlist
\item
  Let H be a connected closed subgroup of G, of maximum rank, distinct from G and maximal for these properties. Denote by Z the quotient group \(\mathrm{C}(\mathrm{H})/\mathrm{C}(\mathrm{G})\) .
\end{enumerate}

\begin{enumerate}
\def\labelenumi{\alph{enumi})}
\item
  We have \(\dim Z \leq 1\) ; if \(\dim Z = 0\) , then Z is of order 2,3 or 5 (reduce to the case in which G is almost simple, with trivial centre; apply Chap. VI, §4, Exerc. 4).
\item
  Assume that G is almost simple and that Z is of order 3 or 5; put \(z = \text{Card}(Z)\) and \(\pi = \text{Card}(\pi_1(H))/\text{Card}(\pi_1(G))\) . Show that we must be in one of the following seven cases:
\end{enumerate}

\begin{enumerate}
\def\labelenumi{(\roman{enumi})}
\item
  G is of type \(\mathrm{G}_2\) , H is of type \(\mathrm{A}_2\) , and \(z = 3, \pi = 1\) ;
\item
  G is of type \(F_{4}\) , H is of type \(A_{2} \times A_{2}\) , and \(z = \pi = 3\) ;
\item
  G is of type \(E_{6}\) , H is of type \(A_{2} \times A_{2} \times A_{2}\) , and \(z = \pi = 3\) ;
\item
  G is of type \(E_{7}\) , H is of type \(A_{2} \times A_{5}\) , and \(z = \pi = 3\) ;
\item
  G is of type \(E_{8}\) , H is of type \(A_{8}\) , and \(z = \pi = 3\) ;
\item
  G is of type \(E_{8}\) , H is of type \(A_{2} \times E_{6}\) , and \(z = \pi = 3\) ;
\item
  G is of type \(E_{8}\) , H is of type \(A_{4} \times A_{4}\) , and \(z = \pi = 5\) .
\end{enumerate}

(Use loc. cit. and the plates of Chap. VI; to calculate \(\pi\) , remark that if \(f\) and \(f'\) denote the connection indices of G and H, respectively, then \(z\pi f = f'\) .)

\begin{enumerate}
\def\labelenumi{\alph{enumi})}
\setcounter{enumi}{2}
\tightlist
\item
  In each of the preceding cases, determine the group H.
\end{enumerate}

\begin{enumerate}
\def\labelenumi{\arabic{enumi})}
\setcounter{enumi}{9}
\tightlist
\item
  We retain the notations of the preceding exercise.
\end{enumerate}

\begin{enumerate}
\def\labelenumi{\alph{enumi})}
\item
  Assume that \(\dim Z = 1\) . Then there are exactly two complex structures on G/H invariant under G; there exists an automorphism of G that leaves H stable and interchanges these two structures (use Exerc. 8).
\item
  Determine the complex structures on \(\mathbf{P}_n(\mathbf{C})\) invariant under \(\mathbf{SU}(n + 1,\mathbf{C})\) .
\item
  Assume that \(\dim Z = 0\) and \(\operatorname{Card}(Z) \neq 2\) . Show that there exist almost complex structures on G/H invariant under G (if z is an element of C(H) that is not central in G, Int z induces an automorphism of G/H of odd order (Exerc. 9); use Exerc. 8 b)). Show that there is no complex structure on G/H invariant under G (use Exerc. 8 d)).
\end{enumerate}

\(d\) ) There is no complex structure on \(\mathbf{S}_6\) invariant under \(\mathbf{SO}(7,\mathbf{R})\) (use Exerc. 7 of §3).

\begin{enumerate}
\def\labelenumi{\arabic{enumi})}
\setcounter{enumi}{10}
\tightlist
\item
  We retain the notations of Exerc. 9 and 10.
\end{enumerate}

\begin{enumerate}
\def\labelenumi{\alph{enumi})}
\item
  The homogeneous space G/H is symmetric (§1, Exerc. 8) if and only if Z is of order 2 or of dimension 1. The symmetric space G/H is then irreducible (loc. cit.).
\item
  Assume that \(\dim Z = 1\) ; denote by \(X\) the complex-analytic manifold \(G / H\) . Show that we are then in one of the following situations:
\end{enumerate}

\begin{enumerate}
\def\labelenumi{(\roman{enumi})}
\item
  \(\mathbf{G}\) is of type \(\mathrm{A}_l\) and \(\mathrm{D(H)}\) of type \(\mathrm{A}_{p - 1}\times \mathrm{A}_{l - p}\) ; \(\mathbf{X}\) is isomorphic to the grassmannian \(\mathbf{G}_p(\mathbf{C}^{l + 1})\) .
\item
  G is of type \(B_{l}\) and D(H) of type \(B_{l-1}\) ; X is isomorphic to the submanifold of \(\mathbf{P}_{2l}(\mathbf{C})\) defined by the vanishing of a separating quadratic form (smooth projective quadric).
\end{enumerate}

(ii') G is of type \(D_{l}\) and \(D(H)\) of type \(D_{l-1}\) ; X is isomorphic to a smooth projective quadric in \(\mathbf{P}_{2l-1}(\mathbf{C})\) .

\begin{enumerate}
\def\labelenumi{(\roman{enumi})}
\setcounter{enumi}{2}
\item
  G is of type \(C_{l}\) and D(H) of type \(A_{l-1}\) ; X is isomorphic to the submanifold of \(\mathbf{G}_{l}(\mathbf{C}^{2l})\) consisting of the maximal isotropic subspaces for the usual alternating bilinear form on \(C^{2l}\) .
\item
  G is of type \(D_{l}\) and \(\mathrm{D}(\mathrm{H})\) of type \(A_{l-1}\) ; X is isomorphic to the submanifold of \(\mathbf{G}_{l}(\mathbf{C}^{2l})\) consisting of the maximal isotropic subspaces for the usual symmetric bilinear form on \(C^{2l}\) .
\item
  G is of type \(E_{6}\) and D(H) of type \(D_{5}\) .
\item
  G is of type \(E_{7}\) and D(H) of type \(E_{6}\) .
\end{enumerate}

\begin{enumerate}
\def\labelenumi{\alph{enumi})}
\setcounter{enumi}{2}
\tightlist
\item
  Assume that \(\operatorname{Card}(Z)=2\) ; give the list of possible situations. If G is of type \(A_{l}\) , \(B_{l}\) , \(C_{l}\) or \(D_{l}\) , show that the real manifold G/H is isomorphic to a grassmannian manifold \(\mathrm{G}_{p}(\mathrm{K}^{q})\) , with K=R, C or H.
\end{enumerate}

¶ 12) Assume that G is simply-connected; for all \(\alpha \in \mathrm{R}(\mathrm{G},\mathrm{T})\) , put \(t_{\alpha} = \exp (\frac{1}{2} K_{\alpha})\) . Put \(\mathrm{N} = \mathrm{N}_{\mathrm{G}}(\mathrm{T})\) , and denote by \(\varphi : \mathrm{N} \to \mathrm{W}\) the canonical map. Let B be a basis of \(\mathrm{R}(\mathrm{G},\mathrm{T})\) ; choose an element \(n_{\alpha}\) of \((\mathrm{N} \cap \mathrm{S}_{\alpha}) - (\mathrm{T} \cap \mathrm{S}_{\alpha})\) for all \(\alpha \in \mathrm{B}\) .

\begin{enumerate}
\def\labelenumi{\alph{enumi})}
\item
  Show that \(\varphi(n_{\alpha}) = s_{\alpha}\) and \(n_{\alpha}^{2} = t_{\alpha}\) , so that \(n_{\alpha}^{4} = 1\) .
\item
  Let \(\alpha\) and \(\beta\) be two distinct elements of B, and let \(m_{\alpha \beta}\) be the order of \(s_{\alpha}s_{\beta}\) in W. Prove that
\end{enumerate}

\[
\begin{array}{r l} {n _ {\alpha} n _ {\beta} = n _ {\beta} n _ {\alpha}} & {\mathrm{if} m _ {\alpha \beta} = 2} \\ {n _ {\alpha} n _ {\beta} n _ {\alpha} = n _ {\beta} n _ {\alpha} n _ {\beta}} & {\mathrm{if} m _ {\alpha \beta} = 3} \\ {(n _ {\alpha} n _ {\beta}) ^ {2} = (n _ {\beta} n _ {\alpha}) ^ {2}} & {\mathrm{if} m _ {\alpha \beta} = 4} \\ {(n _ {\alpha} n _ {\beta}) ^ {3} = (n _ {\beta} n _ {\alpha}) ^ {3}} & {\mathrm{if} m _ {\alpha \beta} = 6} \end{array}
\]

(if for example \(m_{\alpha\beta}=3\) , then \((s_{\alpha}s_{\beta})s_{\alpha}(s_{\alpha}s_{\beta})^{-1}=s_{\beta}\) and \(s_{\alpha}s_{\beta}(\alpha)=\beta\) ; show that \(n_{\alpha}n_{\beta}n_{\alpha}n_{\beta}^{-1}n_{\alpha}^{-1}n_{\beta}^{-1}\) belongs to \(S_{\beta}\) , and conclude by remarking that \(S_{\alpha}\cap S_{\beta}\cap T=\{e\}\) ).

\begin{enumerate}
\def\labelenumi{\alph{enumi})}
\setcounter{enumi}{2}
\item
  Deduce from b) that there exists a unique section \(\nu: W \to N\) of \(\varphi\) such that \(\nu(s_{\alpha}) = n_{\alpha}\) and \(\nu(ww') = \nu(w)\nu(w')\) if \(l(ww') = l(w) + l(w')\) (where \(l(w)\) denotes the length of w with respect to the generating system \((s_{\alpha})_{\alpha \in \mathbb{B}}\) ; cf.~Chap. VI, §1, no. 5, Prop. 5). Put \(n_w = \nu(w)\) .
\item
  Let \(W^{*}\) be the subgroup of N generated by the \(n_{\alpha}\) ; show that \(W^{*} \cap T\) is the subgroup \(T_{2}\) of T consisting of the elements of order \(\leq 2\) and that W can be identified with \(W^{*}/T_{2}\) .
\item
  Let \(w \in \mathrm{W}\) be such that \(w^2 = 1\) ; show that \(n_w^2 = \prod_{\alpha \in \mathrm{R}_w} t_\alpha\) , where \(\mathrm{R}_w\) is the set of positive roots \(\alpha\) such that \(w(\alpha) < 0\) (write \(w = s_{\alpha_1} \ldots s_{\alpha_r}\) , with \(r = l(w)\) and \(\alpha_i \in \mathrm{B}\) ; apply c) and Chap. VI, §1, no. 6, Cor. 2 of Prop. 17).
\end{enumerate}

\(f\) ) Assume that \(\mathrm{G}\) is almost simple. Let \(c\) be a Coxeter transformation in \(\mathrm{W}\) and \(h\) the Coxeter number of \(\mathrm{W}\) (Chap. VI, §1, no. 11). Show that \(n_c^h = \prod_{\alpha \in \mathbb{R}_+}t_\alpha\) (use \(c\) ), \(e\) ) and Exerc. 2 of Chap. V, §6).

\begin{enumerate}
\def\labelenumi{\arabic{enumi})}
\setcounter{enumi}{12}
\tightlist
\item
  \begin{enumerate}
  \def\labelenumii{\alph{enumii})}
  \tightlist
  \item
    Choose a basis B of R(G, T), and denote by R+ the set of positive roots of R(G, T). Prove that \(z_{G} = \prod_{\alpha \in R_{+}} \exp(\frac{1}{2} K_{\alpha})\) is an element of C(G), independent of the choice of T and of B; we have \(z_{G}^{2} = e\) .
  \end{enumerate}
\end{enumerate}

\begin{enumerate}
\def\labelenumi{\alph{enumi})}
\setcounter{enumi}{1}
\item
  Let H be another connected compact Lie group. Then \(z_{\mathrm{G}\times\mathrm{H}} = (z_{\mathrm{G}}, z_{\mathrm{H}})\) ; prove that, if \(f : G \to H\) is a surjective morphism of Lie groups, then \(f(z_{\mathrm{G}}) = z_{\mathrm{H}}\) .
\item
  Put \(\mathrm{R} = \mathrm{R}(\mathrm{G},\mathrm{T})\) ; assume that \(\mathbf{G}\) is simply-connected, so that \(\mathrm{X(T)}\) can be identified with \(\mathrm{P(R)}\) . Show that the kernel of the homomorphism \(\chi \mapsto \chi (z_{\mathrm{G}})\)
\end{enumerate}

from X(T) to \(\{1,-1\}\) is the subgroup \(P'(R)\) defined in Exerc. 8 of Chap. VI, §1.

\begin{enumerate}
\def\labelenumi{\alph{enumi})}
\setcounter{enumi}{3}
\tightlist
\item
  If \(\mathrm{G} = \mathbf{S}\mathbf{U}(n,\mathbf{C})\) , then \(z_{\mathrm{G}} = (-1)^{n + 1}I_n\) ;
\end{enumerate}

if \(\mathrm{G} = \mathbf{S}\mathbf{U}(n,\mathbf{H})\) , then \(z_{\mathrm{G}} = -I_n\) ;

if \(\mathrm{G} = \mathbf{Spin}(n, \mathbf{R})\) with \(n \equiv 3, 4, 5\) or 6 (mod 8), then \(z_{G} = -1\) ;

if \(\mathrm{G} = \mathbf{Spin}(n, \mathbf{R})\) with \(n \equiv 0, 1, 2, 7 \pmod{8}\) , then \(z_{G} = 1\) ;

if G is of type \(E_{6}\) , \(E_{8}\) , \(F_{4}\) or \(G_{2}\) , then \(z_{G} = e\) ;

if G is simply-connected of type \(E_{7}\) , then \(z_{G}\) is the unique non-identity element of \(\mathrm{C}(\mathrm{G})\) .

(Use Chap. VI, §4, Exerc. 5.)

\begin{enumerate}
\def\labelenumi{\arabic{enumi})}
\setcounter{enumi}{13}
\tightlist
\item
  Assume that the group G is almost simple; denote the Coxeter number of R(G,T) by h (Chap. VI, §1, no. 11). An element g of G is said to be a Coxeter element if there exists a maximal torus S of G such that g belongs to \(\mathrm{N}_{\mathrm{G}}(\mathrm{S})\) and its class in \(\mathrm{W}_{\mathrm{G}}(\mathrm{S})\) is a Coxeter transformation (loc. cit.).
\end{enumerate}

\begin{enumerate}
\def\labelenumi{\alph{enumi})}
\item
  Show that any two Coxeter elements are conjugate (argue as in the proof of the Cor. of Prop. 10, no. 5).
\item
  Every Coxeter element g is regular and satisfies \(g^{h} = z_{G}\) , where \(z_{G}\) is the element of C(G) defined in Exerc. 13; in particular, g is of order h or 2h according as \(z_{G}\) is or is not equal to e (use Exerc. 12 f)).
\item
  An element \(g \in G\) is a Coxeter element if and only if the automorphism \(Ad g \otimes 1_{C}\) of \(g_{C}\) satisfies the equivalent conditions of Chap. VIII, §5, Exerc. 5 f).
\end{enumerate}

\(d\) ) Show that every regular element \(g\) of \(\mathbf{G}\) such that \(g^{h}\in \mathrm{C}(\mathrm{G})\) is a Coxeter element; for \(p < h\) , there is no regular element \(k\) such that \(k^p\in \mathrm{C}(\mathrm{G})\) .

\begin{enumerate}
\def\labelenumi{\arabic{enumi})}
\setcounter{enumi}{14}
\tightlist
\item
  Let H be a connected closed subgroup of G. Then H is said to be clean if it is not contained in any connected closed subgroup of maximal rank distinct from G.
\end{enumerate}

\begin{enumerate}
\def\labelenumi{\alph{enumi})}
\item
  Show that H is clean if and only if its centralizer in G is equal to C(G). In particular, if H is clean then C(H) = C(G) ∩ H.
\item
  Assume from now on that H is clean. Show that, for every maximal torus S of H, we have \(\mathrm{C}(\mathrm{H}) = \mathrm{S} \cap \mathrm{C}(\mathrm{G})\) .
\item
  Let \(H'\) be a connected closed subgroup of G containing H. Prove that \(rkH < rkH'\) , and deduce that H is clean in \(H'\) .
\item
  Let K be a connected closed subgroup of H, clean in H and containing a regular element of G. Show that K is clean in G.
\end{enumerate}

\begin{enumerate}
\def\labelenumi{\arabic{enumi})}
\setcounter{enumi}{15}
\tightlist
\item
  Let H be a connected closed subgroup of G such that \(T \cap H\) is a maximal torus S of H. Let \(\lambda \in \mathrm{R}(\mathrm{H}, \mathrm{S})\) ; denote by \(\mathrm{R}(\lambda)\) the set of roots of \(\mathrm{R}(\mathrm{G}, \mathrm{T})\) whose restriction to S is equal to \(\lambda\) .
\end{enumerate}

\begin{enumerate}
\def\labelenumi{\alph{enumi})}
\item
  Show that \(\mathrm{R}(\lambda)\) is not empty.
\item
  Let \(w \in \mathrm{W}_{\mathrm{H}}(\mathrm{S})\) ; show that there exists an element \(\bar{w}\) of \(\mathrm{W}_{\mathrm{G}}(\mathrm{T})\) such that \(\mathrm{R}(w\lambda) = \bar{w}\mathrm{R}(\lambda)\) (use Exerc. 4 of §2).
\item
  Let \(\mathrm{P}(\lambda)\) be the intersection of \(\mathrm{R}(\mathrm{G},\mathrm{T})\) with the subgroup of \(\mathrm{X}(\mathrm{T})\) generated by \(\mathrm{R}(\lambda)\) . Show that there exists a closed subgroup \(G_{\lambda}\) of G containing T such that \(\mathrm{P}(\lambda)=\mathrm{R}(\mathrm{G}_{\lambda},\mathrm{T})\) . Deduce that the reflection \(s_{\lambda}\in\mathrm{W}_{\mathrm{H}}(\mathrm{S})\) is the restriction to S of a product of reflections \(s_{\alpha}\) with \(\alpha\in\mathrm{P}(\lambda)\) .
\item
  Show that the nodal vector \(K_{\lambda} \in \mathrm{L}(\mathrm{S})\) associated to \(\lambda\) is a linear combination with integer coefficients of the \(K_{\alpha}\) for \(\alpha \in \mathrm{R}(\lambda)\) .
\item
  Let \(B_{H}\) be a basis of \(R(H,S)\) . Show that \(R(\lambda)\) is contained in the subgroup of \(X(T)\) generated by the union of the \(R(\mu)\) for \(\mu \in B_{H}\) (prove that the set of \(\lambda \in R(H,S)\) having the stated property is stable under \(s_{\mu}\) for all \(\mu \in B_{H}\) , and use c)).
\end{enumerate}

¶ 17) We retain the notations of the preceding exercise; assume further that the subgroup H is clean (Exerc. 15).

\begin{enumerate}
\def\labelenumi{\alph{enumi})}
\item
  Show that \(\mathrm{R}(\mathrm{G},\mathrm{T})\) is contained in the subgroup of \(\mathrm{X}(\mathrm{T})\) generated by the union of the \(\mathrm{R}(\lambda)\) for \(\lambda\in B_{H}\) .
\item
  Let \(\Delta\) be the identity component of the subgroup of S consisting of the \(s \in S\) such that \(\lambda(s) = \mu(s)\) for all \(\lambda, \mu\) in \(B_{H}\) . Show that there exists a basis \(\{\alpha_{1}, \ldots, \alpha_{l}\}\) of R(G,T) and an integer k, with \(0 \leq k \leq l - 1\) , such that \(\Delta\) is the identity component of the set of \(t \in T\) satisfying
\end{enumerate}

\[
\alpha_ {1} (t) = \dots = \alpha_ {k} (t) = 1, \alpha_ {k + 1} (t) = \dots = \alpha_ {l} (t).
\]

(Let \(x\) be the element of \(\mathrm{L}(\mathrm{S})\) such that \(\delta(\lambda)(x) = 2\pi i\) for all \(\lambda \in \mathrm{B}_{\mathrm{H}}\) ; deduce from \(a\) ) that \(\exp x \in \mathrm{C}(\mathrm{G})\) . Choose \(\{\alpha_1, \ldots, \alpha_l\}\) so that \(i\delta(\alpha_j)(x)\) is zero for \(1 \leq j \leq k\) and \(< 0\) for \(k + 1 \leq j \leq l\) ; then show that for all \(\lambda \in \mathrm{B}_{\mathrm{H}}\) every root in \(\mathrm{R}(\lambda)\) can be written

\[
\alpha_ {j} + n _ {1} \alpha_ {1} + \dots + n _ {k} \alpha_ {k},
\]

with \(j \geq k + 1\) and \(n_i \in \mathbf{N}\) . Conclude by using \(a)\) .

\begin{enumerate}
\def\labelenumi{\alph{enumi})}
\setcounter{enumi}{2}
\tightlist
\item
  Show that the integer k does not depend on the choice of the tori S, T or of the bases \(B_{H}, \{\alpha_{1}, \ldots, \alpha_{l}\}\) ; it is equal to the rank of the derived group of \(Z_{\mathrm{G}}(\Delta)\) .
\end{enumerate}

\begin{enumerate}
\def\labelenumi{\arabic{enumi})}
\setcounter{enumi}{17}
\tightlist
\item
  We retain the notations of Exerc. 16 and 17. The subgroup H is said to be principal if it is clean and if the sub-torus \(\Delta\) contains a regular element (in other words, if k = 0).
\end{enumerate}

\begin{enumerate}
\def\labelenumi{\alph{enumi})}
\item
  Let \(\mathrm{K}\) be a connected closed subgroup of \(\mathrm{H}\) . Show that \(\mathrm{K}\) is a principal subgroup of \(\mathrm{G}\) if and only if \(\mathrm{K}\) is principal in \(\mathrm{H}\) and \(\mathrm{H}\) is principal in \(\mathrm{G}\) (use Exerc. 15 c) and d)).
\item
  Assume from now on that H is principal. Show that the union of the \(\mathrm{R}(\mu)\) , for \(\mu \in B_{H}\) , is a basis of \(\mathrm{R}(\mathrm{G}, \mathrm{T})\) .
\item
  Let \(\alpha \in \mathrm{R}(\mathrm{G},\mathrm{T})\) . Prove that the restriction of \(\alpha\) to S is a non-zero integer multiple of a root \(\lambda \in \mathrm{R}(\mathrm{H},\mathrm{S})\) ; the root \(\alpha\) is a sum of elements of \(\mathrm{R}(\lambda)\) (show that the intersection of \(\mathrm{L}(\mathrm{S})\) with the hyperplane \(\delta (\alpha) = 0\) is a wall of \(\mathrm{L}(\mathrm{S}))\) ).
\item
  Let \(\lambda\in\mathrm{R}(\mathrm{H},\mathrm{S})\) ; show that the nodal vector \(K_{\lambda}\) is a linear combination with non-negative integer coefficients of the \(K_{\alpha}\) for \(\alpha\in\mathrm{R}(\lambda)\) (cf.~Exerc. 16 d) and Chap. V, §3, no. 5, Lemma 6).
\end{enumerate}

\begin{enumerate}
\def\labelenumi{\arabic{enumi})}
\setcounter{enumi}{18}
\tightlist
\item
  We retain the notations of Exerc. 16 to 18; assume that the subgroup H is principal. Give \(\mathrm{X(T)\otimes\mathbf{R}}\) (resp. \(\mathrm{X(S)\otimes\mathbf{R}}\) ) a scalar product invariant under \(\mathrm{W_{G}(T)}\) (resp. under \(\mathrm{W_{H}(S)}\) ). Let \(\lambda,\mu\) be two roots in \(B_{H}\) .
\end{enumerate}

\begin{enumerate}
\def\labelenumi{\alph{enumi})}
\item
  Show that, if \(\lambda\) and \(\mu\) are orthogonal, the sets \(\mathrm{R}(\lambda)\) and \(\mathrm{R}(\mu)\) are orthogonal. Deduce that if \(\mathrm{G}\) is almost simple then so is \(\mathrm{H}\) .
\item
  Assume from now on that \(n(\lambda, \mu) = -1\) . Show that there exists a surjective map \(u: \mathrm{R}(\lambda) \to \mathrm{R}(\mu)\) such that, for \(\alpha \in \mathrm{R}(\lambda)\) , \(u(\alpha)\) is the unique root of \(\mathrm{R}(\mu)\) linked (that is, not orthogonal) to \(\alpha\) ; we have \(K_{\mu} = \sum_{\beta \in \mathrm{R}(\mu)} K_{\beta}\) , and the roots in \(\mathrm{R}(\mu)\) are pairwise orthogonal (write \(K_{\mu}\) and \(K_{\lambda}\) as linear combinations of the \(K_{\alpha}\) , then identify the coefficients).
\item
  If \(n(\mu, \lambda) = -1\) , the map \(u\) is bijective, and the only pairs of linked roots in \(\mathrm{R}(\lambda) \cup \mathrm{R}(\mu)\) are the pairs \((\alpha, u(\alpha))\) for \(\alpha \in \mathrm{R}(\lambda)\) .
\item
  Assume that \(n(\mu,\lambda)=-2\) ; let \(\beta\in\mathrm{R}(\mu)\) . Show that \(u^{-1}(\beta)\) contains either one or two elements; if \(u^{-1}(\beta)=\{\alpha_{1},\alpha_{2}\}\) , the root \(\alpha_{i}\) (i=1,2) is orthogonal to the other roots in \(\mathrm{R}(\lambda)\) , and has the same length as \(\beta\) ; if \(u^{-1}(\beta)=\{\alpha\}\) and \(\|\alpha\|=\|\beta\|\) , \(\alpha\) is linked to a unique root \(\alpha^{\prime}\in\mathrm{R}(\lambda)\) , of the same length as \(\alpha\) , and \(\{\alpha,\alpha^{\prime}\}\) is orthogonal to the remainder of \(\mathrm{R}(\lambda)\) ; if \(u^{-1}(\beta)=\{\alpha\}\) and \(\|\alpha\|\neq\|\beta\|\) , \(\alpha\) is orthogonal to the other roots in \(\mathrm{R}(\lambda)\) .
\item
  Study similarly the case \(n(\mu, \lambda) = -3\) .
\end{enumerate}

\begin{enumerate}
\def\labelenumi{\arabic{enumi})}
\setcounter{enumi}{19}
\tightlist
\item
  Assume that the group G is almost simple; let H be a principal subgroup of G, of rank \(\geq 2\) .
\end{enumerate}

\begin{enumerate}
\def\labelenumi{\alph{enumi})}
\item
  Show that H is semi-simple of type \(B_{h}, C_{h}, F_{4}\) or \(G_{2}\) (use Exerc. 19).
\item
  Assume that H is of rank \(\geq 3\) . Show that G is of type \(A_{l}, D_{l}, E_{6}, E_{7}\) or \(E_{8}\) (consider the terminal vertices of the Dynkin graph of R(G, T), and apply Exerc. 19).
\item
  Show that, if \(rkH \geq 3\) , we are in one of the following situations:
\end{enumerate}

G is of type \(A_{2l}\) \((l \geq 3)\) and H of type \(B_{l}\) ;

G is of type \(\mathrm{A}_{2l - 1}\) \((l\geq 3)\) and H of type \(\mathbf{C}_l\)

G is of type \(D_{l}\) \((l \geq 4)\) and H of type \(B_{l-1}\) ;

G is of type \(\mathrm{E}_6\) and H of type \(\mathrm{F}_4\) .

\(d\) ) Show that, if H is of type \(\mathrm{B}_2\) , then G is of type \(\mathrm{A}_3\) or \(\mathrm{A}_4\) .

\begin{enumerate}
\def\labelenumi{\alph{enumi})}
\setcounter{enumi}{4}
\tightlist
\item
  Show that, if H is of type \(G_{2}\) , then G is of type \(B_{3}, D_{4}\) or \(A_{6}\) .
\end{enumerate}

(For a more explicit description of these situations, see §5, Exerc. 5.)

\(f)\) Let \(\mathrm{K}\) be a connected closed subgroup of \(\mathrm{G}\) containing \(\mathrm{H}\) ; show that either \(\mathrm{K} = \mathrm{G}\) or \(\mathrm{K} = \mathrm{H}\) , or \(\mathrm{H}\) is of type \(\mathrm{G}_2\) , \(\mathrm{K}\) of type \(\mathrm{B}_3\) and \(\mathrm{G}\) of type \(\mathrm{D}_4\) or \(\mathrm{A}_6\) .

\begin{enumerate}
\def\labelenumi{\arabic{enumi})}
\setcounter{enumi}{20}
\tightlist
\item
  \begin{enumerate}
  \def\labelenumii{\alph{enumii})}
  \tightlist
  \item
    Let H be a connected closed subgroup of G of rank 1. Show that the following conditions are equivalent:
  \end{enumerate}
\end{enumerate}

\begin{enumerate}
\def\labelenumi{(\roman{enumi})}
\item
  H is principal;
\item
  H is clean and contains a regular element of G;
\item
  there exists a principal \(\mathfrak{sl}_2\) -triplet \((x,h,y)\) in \(\mathfrak{g}_{\mathbf{C}}\) (Chap. VIII, §11, no. 4) such that
\end{enumerate}

\[
\mathrm{L} (\mathrm{H}) _ {(\mathbf {C})} = \mathbf {C} x + \mathbf {C} h + \mathbf {C} y.
\]

\begin{enumerate}
\def\labelenumi{\alph{enumi})}
\setcounter{enumi}{1}
\item
  Show that G contains a principal subgroup of rank 1, and that any two such subgroups are conjugate (with the notations of Exerc. 17, remark that \(S = \Delta\) ).
\item
  Show that a connected closed subgroup of G is principal if and only if it contains a principal subgroup of G of rank 1 (use Exerc. 18 a)).
\end{enumerate}

\begin{enumerate}
\def\labelenumi{\arabic{enumi})}
\setcounter{enumi}{21}
\item
  Let H be a principal subgroup of G of rank 1; let \(\Gamma\) be the subgroup of \(\operatorname{Aut}(G)\) consisting of the automorphisms u such that \(u(\mathrm{H}) = \mathrm{H}\) . Show that \(\operatorname{Aut}(G) = \Gamma.\operatorname{Int}(G)\) .
\end{enumerate}

\exercisesection{\S~5}

\begin{enumerate}
\def\labelenumi{\arabic{enumi})}
\item
  Assume that G is simply-connected; the alcoves of G are the subsets of G of the form \(\exp(A)\) , where A is an alcove of a Cartan subalgebra of g.

\begin{enumerate}
\def\labelenumi{\alph{enumi})}
\item
  Show that the alcoves of \(\mathbf{G}\) form a partition of \(\mathrm{G}_r\) .
\item
  Every alcove of G is contained in a unique maximal torus.
\item
  Show that the alcoves of G contained in T form a partition of \(T_{r}\) , and that the set of such alcoves is a principal homogeneous space under W.
\item
  Every conjugacy class of regular elements of G has a unique point in each alcove.
\item
  Let E be an alcove of G. Show that E is a contractible space; if g is simple, E is homeomorphic to an open simplex in a euclidean space.
\end{enumerate}

\begin{enumerate}
\def\labelenumi{\arabic{enumi})}
\setcounter{enumi}{1}
\tightlist
\item
  \begin{enumerate}
  \def\labelenumii{\alph{enumii})}
  \tightlist
  \item
    Let E be a simply-connected topological space and \(u : E \to G\) a continuous map such that \(u(E) \subset G_r\) . Show that u is homotopic (General Topology, Chap. XI) to the constant map with value e (consider the covering \(\varphi_r : (G/T) \times t_r \to G_r\) ; lift u to \(\tilde{u} : E \to (G/T) \times t\) , then use the fact that t is contractible).
  \end{enumerate}
\end{enumerate}

\(^{*}b)\) Prove that the group \(\pi_{2}(G)\) is zero.

(Let \(u: \mathbf{S}_2 \to \mathbf{G}\) be a map of class \(\mathrm{C}^\infty\) ; by using Prop. 1 and the transversality theorem, show that \(u\) is homotopic to a map with image contained in \(\mathrm{G}_r\) , then apply \(a).)_*\)

\begin{enumerate}
\def\labelenumi{\arabic{enumi})}
\setcounter{enumi}{2}
\tightlist
\item
  Let \(f: \tilde{\mathrm{G}} \to \mathrm{G}\) be a universal covering of \(\mathrm{G}\) and \(\pi\) the kernel of \(f\) (isomorphic to \(\pi_1(\mathrm{G})\) ); put \(\mathrm{C} = \mathrm{C}(\mathrm{G})\) and \(\tilde{\mathrm{C}} = \mathrm{C}(\tilde{\mathrm{G}})\) . Let \(\sigma\) be an automorphism of
\end{enumerate}

G, \(\tilde{\sigma}\) the automorphism of \(\tilde{G}\) induced by \(\sigma\) . Denote by \(G_{\sigma}, C_{\sigma}, \tilde{G}_{\sigma}, \tilde{C}_{\sigma}\) the set of points of G, C, \(\tilde{G}\) , \(\tilde{C}\) fixed by \(\sigma, \sigma, \tilde{\sigma}, \tilde{\sigma}\) , respectively.

\begin{enumerate}
\def\labelenumi{\alph{enumi})}
\item
  Show that the identity component \((\mathrm{G}_{\sigma})_0\) of \(\mathrm{G}_{\sigma}\) is \(f(\tilde{\mathrm{G}}_{\sigma})\) .
\item
  Denote by s the endomorphism of the Z-module \(\pi\) induced by \(\tilde{\sigma}\) . Show that the quotient group \(\mathrm{G}_{\sigma}/(\mathrm{G}_{\sigma})_{0}\) is isomorphic to a subgroup of \(\operatorname{Coker}(1-s)\) , and in particular is commutative. If 1-s is surjective, \(G_{\sigma}\) is connected.
\item
  Let \(n \in N\) be such that \(s^{n} = \operatorname{Id}_{\pi}\) . Show that the group \(\mathrm{G}_{\sigma}/(\mathrm{G}_{\sigma})_{0}\) can be identified with a subgroup of the quotient \(\operatorname{Ker}(1 + s + \cdots + s^{n-1})/\operatorname{Im}(1 - s)\) ; deduce that it is annihilated by n.~If n is prime to the order of the torsion subgroup of \(\pi\) , then \(G_{\sigma}\) is connected.
\end{enumerate}

\(d\) ) Recover the results of \(b\) ) and \(c\) ) by using Exerc. 23 of Algèbre, Chap. X, p.~194.

\begin{enumerate}
\def\labelenumi{\alph{enumi})}
\setcounter{enumi}{4}
\tightlist
\item
  Show that \(\mathrm{G}_{\sigma}\) is connected in each of the following cases:
\end{enumerate}

\begin{enumerate}
\def\labelenumi{(\roman{enumi})}
\item
  G is semi-simple of type \(\mathrm{A}_{2n}\) ( \(n \geq 1\) ) and \(\sigma\) is not inner;
\item
  G is semi-simple of type \(\mathrm{E}_6\) and \(\sigma\) is not inner;
\item
  G is semi-simple of type \(\mathrm{D}_4\) and \(\sigma\) is a triality (that is, of order 3 modulo \(\operatorname{Int}(\mathbf{G})\) ).
\end{enumerate}

\(f\) ) Define an isomorphism from \(\mathrm{C}_{\sigma} / (\mathrm{C}_{\sigma}\cap (\mathrm{G}_{\sigma})_{0})\) to \(((1 - \tilde{\sigma})\tilde{\mathrm{C}}\cap \pi) / (1 - \tilde{\sigma})\pi .\) Deduce that if \(1 - s\) is not surjective and \(\pi \subset (1 - \tilde{\sigma})\mathrm{C}\) , then \(\mathrm{C}_{\sigma}\) is not contained in \((\mathrm{G}_{\sigma})_{0}\) . For \(\mathrm{G} = \mathbf{SO}(2n,\mathbf{R})\) and \(\sigma\) not inner, we have \(-I_{2n}\notin (\mathrm{G}_{\sigma})_{0}\) and \(\mathrm{G}_{\sigma}\) is not connected.

\begin{enumerate}
\def\labelenumi{\arabic{enumi})}
\setcounter{enumi}{3}
\tightlist
\item
  Assume that G is semi-simple. Let e be a framing of G and \(\Phi\) a group of automorphisms of G respecting this framing; denote by H the subgroup of G consisting of the elements fixed by \(\Phi\) .
\end{enumerate}

\begin{enumerate}
\def\labelenumi{\alph{enumi})}
\item
  Show that \(H_{0}\) is semi-simple; if G is simply-connected, H is connected (reduce to the case in which G is almost simple and \(\Phi\) is cyclic, and apply Th. 1 (no. 3) and Chap. VIII, §5, Exerc. 13).
\item
  Assume that \(\mathbf{G}\) is almost simple. Show that
\end{enumerate}

if G is of type \(\mathrm{A}_{2l}\) ( \(l \geq 1\) ) and \(\varPhi\) is of order 2, \(\mathrm{H}_0\) is of type \(\mathrm{B}_l\) ;

if G is of type \(A_{2l-1}\) ( \(l \geq 2\) ) and \(\Phi\) is of order 2, \(H_{0}\) is of type \(C_{l}\) ;

if G is of type \(D_{l}\) \((l \geq 4)\) and \(\Phi\) is of order 2, \(H_{0}\) is of type \(B_{l-1}\) ;

if G is of type \(D_{4}\) and \(\Phi\) is of order 3 or 6, \(H_{0}\) is of type \(G_{2}\) ;

if \(\mathrm{G}\) is of type \(\mathrm{E}_6\) and \(\varPhi\) is of order 2, \(\mathrm{H}_0\) is of type \(\mathrm{F}_4\) .

In each case, determine the groups \(\pi_i(\mathrm{H})\) , \(i = 0,1\) (cf.~Exerc. 3).

\begin{enumerate}
\def\labelenumi{\alph{enumi})}
\setcounter{enumi}{2}
\tightlist
\item
  Prove that the subgroup \(\mathrm{H}_0\) is principal (§4, Exerc. 18).
\end{enumerate}

\(d\) ) Prove that a semi-simple group of type \(\mathrm{B}_3\) or \(\mathrm{A}_6\) contains a principal subgroup of type \(\mathrm{G}_2\) (use \(b\) ) and Exerc. 20 of §4).

\begin{enumerate}
\def\labelenumi{\arabic{enumi})}
\setcounter{enumi}{4}
\tightlist
\item
  Assume that the group G is almost simple; let H be a principal connected closed subgroup of G (§4, Exerc. 18). Denote by \(\Phi\) the group of automorphisms u of G that fix H, and by F the subgroup of elements of G fixed by \(\Phi\) .
\end{enumerate}

\begin{enumerate}
\def\labelenumi{\alph{enumi})}
\item
  Show that there exists a framing of \(\mathrm{G}\) stable under \(\varPhi\) .
\item
  Show that we must be in one of the following situations:
\end{enumerate}

\begin{enumerate}
\def\labelenumi{(\roman{enumi})}
\item
  \(H = F_{0};\)
\item
  G is of type \(B_{3}\) , H of type \(G_{2}\) and \(\Phi\) reduces to the identity;
\item
  G is of type \(A_{6}\) , H of type \(G_{2}\) and \(\Phi\) of order 2.
\end{enumerate}

(Use Exerc. 4 and Exerc. 20 of §4.)

¶ 6) Assume that G is simply-connected. Let p be a prime number and g an element of C(G) such that \(g^{p} = e\) .

\begin{enumerate}
\def\labelenumi{\alph{enumi})}
\tightlist
\item
  Show that there exist elements \(u \in T\) and \(w \in W\) such that
\end{enumerate}

\begin{enumerate}
\def\labelenumi{(\roman{enumi})}
\item
  \(w(u)u^{-1} = g;\)
\item
  \(w^{p} = 1;\)
\item
  \(u^{p} = e\) if \(p \neq 2\) , \(u^{p} = e\) or g if p = 2.
\end{enumerate}

(Let A be an alcove in t; for \(i = 0, 1, \ldots, p - 1\) , let \(x_i \in \overline{A}\) be such that \(\exp x_i = g^i\) . Take u to be the element \(\exp x\) , where x is the barycentre of the facet of \(\overline{A}\) whose vertices are the \(x_i\) , and w to be the element of W such that \(w(\mathrm{A}) = \mathrm{A} - x\) .)

\begin{enumerate}
\def\labelenumi{\alph{enumi})}
\setcounter{enumi}{1}
\tightlist
\item
  Prove that there exist \(u \in \mathrm{T}\) and \(v \in \mathrm{N}_{\mathrm{G}}(\mathrm{T})\) such that
\end{enumerate}

\begin{enumerate}
\def\labelenumi{(\roman{enumi})}
\item
  \(vuv^{-1}u^{-1}=g;\)
\item
  \(v^{p} = e\) if \(p \neq 2\) , \(v^{p} = e\) or g if p = 2;
\item
  \(u^{p} = e\) if \(p \neq 2\) , \(u^{p} = e\) or \(g\) if \(p = 2\) .
\end{enumerate}

(Use the construction in \(a\) ), lifting \(w\) to \(n_w\) as in Exerc. 12 of §4; take \(v = n_w^{p+1}\) if \(p \neq 2\) , \(v = n_w\) if \(p = 2\) .)

\begin{enumerate}
\def\labelenumi{\arabic{enumi})}
\setcounter{enumi}{6}
\tightlist
\item
  Let p be a prime number. Show that the following conditions are equivalent: (i) p does not divide the order of the torsion subgroup of \(\pi_{1}(G)\) ;
\end{enumerate}

\begin{enumerate}
\def\labelenumi{(\roman{enumi})}
\setcounter{enumi}{1}
\tightlist
\item
  for every element g of G of order p, the centralizer of g in G is connected; (iii) every subgroup of G isomorphic to \((\mathbf{Z}/p\mathbf{Z})^{2}\) is contained in a maximal torus.
\end{enumerate}

(To prove that (i) \(\Longrightarrow\) (ii), use Exerc. 3 \(c\) ); to prove that (iii) \(\Longrightarrow\) (i), use Exerc. 6.)

\begin{enumerate}
\def\labelenumi{\arabic{enumi})}
\setcounter{enumi}{7}
\tightlist
\item
  Take \(\mathrm{G} = \mathbf{SO}(8, \mathbf{R}) / \{\pm I_8\}\) ; denote by \((\alpha_i)_{1 \leq i \leq 4}\) a basis of \(\mathrm{R}(\mathrm{G}, \mathrm{T})\) such that \(\alpha_1, \alpha_3\) and \(\alpha_4\) are not orthogonal to \(\alpha_2\) (Chap. VI, Plate IV). Let A be the subgroup of T consisting of the \(t \in T\) such that
\end{enumerate}

\[
\alpha_ {1} (t) = \alpha_ {3} (t) = \alpha_ {4} (t), \alpha_ {1} (t) ^ {2} = \alpha_ {2} (t) ^ {2} = 1.
\]

Prove that A is isomorphic to \((\mathbf{Z}/2\mathbf{Z})^{2}\) and that its centralizer in G is a finite non-commutative group.

\begin{enumerate}
\def\labelenumi{\arabic{enumi})}
\setcounter{enumi}{8}
\tightlist
\item
  Let R be an irreducible root system, \(R^{\vee}\) the inverse system, B a basis of R and \(\alpha\) the highest root of R (relative to B). Put
\end{enumerate}

\[
\alpha = \sum_ {\beta \in \mathrm{B}} n _ {\beta} \beta \text { and } \alpha^ {\vee} = \sum_ {\beta \in \mathrm{B}} n _ {\beta} ^ {\vee} \beta^ {\vee};
\]

let \(\nu (\mathrm{R}) = \sup_{\beta \in \mathrm{B}}n_{\beta}^{\vee}\)

\begin{enumerate}
\def\labelenumi{\alph{enumi})}
\item
  Show that the interval \(\mathbf{1},\nu(\mathrm{R})\mathbf{I}\) in N is the union of 1 and the \(n_{\beta}^{\vee}\) for \(\beta\in B\) (put \(\alpha_{0}=\alpha\) ; let \((\alpha_{1},\ldots,\alpha_{q})\) be a sequence of distinct elements of B such that \(\alpha_{i}\) is not orthogonal to \(\alpha_{i-1}\) for \(i=1,\ldots,q\) , such that \(n_{\alpha_{q}}^{\vee}=\nu(\mathrm{R})\) and such that q is maximal with these properties; prove that \(n_{\alpha_{i}}^{\vee}=i+1\) for \(i=1,\ldots,q\) ).
\item
  Let \(p\) be a prime number. Show that the following three properties are equivalent:
\end{enumerate}

\begin{enumerate}
\def\labelenumi{(\roman{enumi})}
\item
  \(p \leq \nu(\mathrm{R});\)
\item
  there exists \(\beta \in \mathrm{B}\) such that \(p = n_{\beta}^{\vee}\) ;
\item
  there exists \(\beta \in \mathrm{B}\) such that \(p\) divides \(n_{\beta}^{\vee}\) .
\end{enumerate}

In this case p is called a torsion prime number of R.

\begin{enumerate}
\def\labelenumi{\alph{enumi})}
\setcounter{enumi}{2}
\tightlist
\item
  For each type of irreducible root system, give the value of \(\nu(\mathrm{R})\) and the torsion prime numbers. Show that the set of the \(n_{\beta}\) and that of the \(n_{\beta}^{\vee}\) coincide except for type \(G_{2}\) .
\end{enumerate}

\(d\) ) Let \(\mathbf{R}'\) be a symmetric closed subset of \(\mathbf{R}\) that is irreducible as a root system. Show that \(\nu (\mathrm{R}^{\prime})\leq \nu (\mathrm{R})\)

\begin{enumerate}
\def\labelenumi{\arabic{enumi})}
\setcounter{enumi}{9}
\tightlist
\item
  Let \(\mathbf{R}\) be a root system and \(p\) a prime number. Show that the following properties are equivalent:
\end{enumerate}

\begin{enumerate}
\def\labelenumi{(\roman{enumi})}
\item
  p is a torsion prime number of an irreducible component of R (Exerc. 9 b));\\
\item
  there exists a symmetric closed subset \(R_{1}\) of R, distinct from R and maximal for these properties, such that \((\mathrm{Q}(\mathrm{R}^{\vee}):\mathrm{Q}(\mathrm{R}_{1}^{\vee}))=p(\mathrm{Q}(\mathrm{R}^{\vee}))\) denotes the Z-module generated by the inverse roots of R and \(\mathrm{Q}(\mathrm{R}_{1}^{\vee})\) the submodule generated by the \(\alpha^{\vee}\) for \(\alpha\in R_{1}\) );
\item
  there exists a symmetric closed subset \(R_{1}\) of R such that the p-torsion submodule of \(\mathrm{Q}(\mathrm{R}^{\vee})/\mathrm{Q}(\mathrm{R}_{1}^{\vee})\) is non-zero.
\end{enumerate}

(To prove that (i) \(\Longrightarrow\) (ii), use Exerc. 4 of Chap. VI, §4; to prove that (iii) \(\Longrightarrow\) (i), use Exerc. 9 d).)

In this case \(p\) is said to be a torsion prime number of \(\mathbb{R}\) .

\begin{enumerate}
\def\labelenumi{\arabic{enumi})}
\setcounter{enumi}{10}
\tightlist
\item
  Let p be a prime number. Then p is said to be a torsion prime number of G if there exists a connected closed subgroup H of G, of maximum rank, such that the p-torsion subgroup of \(\pi_{1}(H)\) is non-zero. Put \(R = R(G, T)\) .
\end{enumerate}

\begin{enumerate}
\def\labelenumi{\alph{enumi})}
\item
  The torsion prime numbers of G are the torsion prime numbers of R (Exerc. 10) and the prime numbers that divide the order of the torsion group of \(\pi_1(\mathrm{G})\) .
\item
  Show that every torsion prime number of G divides \(w/l!\) , where w is the order of W and l the rank of the semi-simple group D(G) (use Chap. VI, §2, no. 4, Prop. 7).
\item
  Let \(t \in T\) and \(n \in N\) be such that \(t^{n} \in \mathrm{C}(\mathrm{G})\) . Let \(R_{1}\) be the set of \(\alpha \in R\) such that \(\alpha(t) = 1\) , and m the order of the torsion group of \(\mathrm{Q}(\mathrm{R}^{\vee}) / \mathrm{Q}(\mathrm{R}_{1}^{\vee})\) (cf.~Exerc. 10). Show that every prime factor of m divides n (reduce to the case in which G is simply-connected and almost simple).
\item
  Assume that G is simply-connected. Let \(g \in G\) , \(n \in N\) be such that \(g^{n}\) belongs to \(\mathrm{C}(\mathrm{G})\) and such that no torsion prime number of G divides n.~Show that the derived group of the centralizer of g is simply-connected.
\item
  Let g be an element of G and p a prime number, such that \(g^{p} = e\) and p is not a torsion prime number of G. Show that the centralizer \(Z(g)\) is connected and that p is not a torsion prime number of \(Z(g)\) .
\item
  Assume that G is simply-connected, and let p be a torsion prime number of G. Prove that there exists an element g of G of order p such that \(\pi_{1}(\mathrm{Z}(g))\) is cyclic of order p (use Exerc. 10 and Exerc. 4 of Chap. VI, §4).
\end{enumerate}

\begin{enumerate}
\def\labelenumi{\arabic{enumi})}
\setcounter{enumi}{11}
\tightlist
\item
  Let \(p\) be a prime number. Show that the following conditions are equivalent:
\end{enumerate}

\begin{enumerate}
\def\labelenumi{(\roman{enumi})}
\item
  \(p\) is not a torsion prime number of \(\mathbf{G}\) ;
\item
  for every subgroup \(\mathbf{F}\) of \(\mathbf{G}\) isomorphic to \((\mathbf{Z} / p\mathbf{Z})^n\) (for \(n\in \mathbf{N}\) ), the centralizer of \(\mathbf{F}\) in \(\mathbf{G}\) is connected;
\end{enumerate}

(ii') for every subgroup F of G isomorphic to \((\mathbf{Z}/p\mathbf{Z})^{2}\) , the centralizer of F in G is connected;

\begin{enumerate}
\def\labelenumi{(\roman{enumi})}
\setcounter{enumi}{2}
\tightlist
\item
  every subgroup of G isomorphic to \((\mathbf{Z}/p\mathbf{Z})^{n}\) for some integer n is contained in a maximal torus;
\end{enumerate}

(iii') every subgroup of G isomorphic to \((\mathbf{Z}/p\mathbf{Z})^{3}\) is contained in a maximal torus.

(To prove that (i) \(\Longrightarrow\) (ii), use Exerc. 11 e); to prove that (iii) \(\Longrightarrow\) (i), use Exerc. 11 f) and Exerc. 7.)

\exercisesection{\S~6}

\begin{enumerate}
\def\labelenumi{\arabic{enumi})}
\tightlist
\item
  Let R be a reduced root system in a real vector space V. Give the space V a scalar product invariant under W(R), and the space S(V) the corresponding scalar product (Topological Vector Spaces, Chap. V, §3, no. 3, formulas (13) and (14)); thus, for \(x_{1},\ldots,x_{n},y_{1},\ldots,y_{n}\) in V,
\end{enumerate}

\[
(x _ {1} \dots x _ {n} | y _ {1} \dots y _ {n}) = \sum_ {\sigma \in \mathfrak {S} _ {n}} (x _ {1} | y _ {\sigma (1)}) \dots (x _ {n} | y _ {\sigma (n)}).
\]

Choose a chamber C of R; put \(N = \text{Card}(R_{+})\) , \(\rho = \frac{1}{2} \sum_{\alpha \in R_{+}} \alpha\) and \(w(R) = \text{Card}(W(R))\) . Let P be the element \(\prod_{\alpha \in R_{+}} \alpha\) of \(\mathbf{S}^{\mathrm{N}}(\mathrm{V})\) .

\(a)\) Show that \(\mathrm{P} = \frac{1}{\mathrm{N!}}\sum_{w\in \mathrm{W(R)}}\varepsilon (w)(w\rho)^{\mathrm{N}}\) (cf.~Chap. VI, §3, no. 3, Prop. 2).

\begin{enumerate}
\def\labelenumi{\alph{enumi})}
\setcounter{enumi}{1}
\item
  Deduce the equality \((\mathrm{P}|\mathrm{P}) = w(\mathrm{R})\prod_{\alpha \in \mathrm{R}_{+}}(\rho |\alpha)\) .
\item
  Prove that \((\mathrm{P}|\mathrm{P}) = 2^{-\mathrm{N}}w(\mathrm{R})\prod_{i = 1}^{l}m_i!\prod_{\alpha \in \mathrm{R}_+}(\alpha |\alpha) = 2^{-\mathrm{N}}\prod_{i = 1}^{l}(m_i + 1)!\prod_{\alpha \in \mathrm{R}_+}(\alpha |\alpha)\) , where \(m_{1},\ldots ,m_{l}\) are the exponents of \(\mathrm{W(R)}\) (use Exerc. 3 of Chap. VIII, §9).
\end{enumerate}

¶ 2) Let V be a finite dimensional real Hilbert space over R. Give the space \(\mathbf{S}(\mathrm{V}^{*})\) of polynomials on V the scalar product defined in Topological Vector Spaces, Chap. V, §3, no. 3, formulas (13) and (14) (cf.~Exerc. 1). Denote by \(\gamma\) the canonical gaussian measure on V (Integration, Chap. IX, §6, nos. 4 to 6); thus, if \((x_{i})_{1\leq i\leq n}\) is an orthonormal basis of \(V^{*}\) , we have \(d\gamma(x)=(2\pi)^{-n/2}e^{-(x|x)/2}dx_{1}\ldots dx_{n}\) .

\begin{enumerate}
\def\labelenumi{\alph{enumi})}
\tightlist
\item
  Let \(q\) be the element of \(\mathbf{S}^2(\mathrm{V})\) that defines the scalar product on \(\mathrm{V}^*\) , and let \(\Delta: \mathbf{S}(\mathrm{V}^*) \to \mathbf{S}(\mathrm{V}^*)\) be the operator given by taking the inner product with \(q\) (Algebra, Chap. III, §11, no. 9). Show that, for every orthonormal basis \((x_i)_{1 \leq i \leq n}\) of \(\mathrm{V}^*\) and every polynomial \(\mathrm{P} \in \mathbf{S}(\mathrm{V}^*)\) , we have \(\Delta(\mathrm{P}) = \frac{1}{2} \sum_{i=1}^{n} \frac{\partial^2\mathrm{P}}{\partial x_i^2}\) . b) For \(\mathrm{P} \in \mathbf{S}(\mathrm{V}^*)\) , put \(\mathrm{P}^* = \mathrm{P} * \gamma\) , so that \(\mathrm{P}^*(x) = \int_{\mathrm{V}} \mathrm{P}(x - y) d\gamma(y)\) for \(x \in \mathrm{V}\) . Prove the equality \(\mathrm{P}^* = \sum_{n=1}^{\infty} \frac{\Delta^n(\mathrm{P})}{n!} = e^{\Delta}(\mathrm{P})\) .
\end{enumerate}

(Reduce to proving the analogous formula for the function \(x \mapsto e^{i(x|u)}\) , for \(u \in \mathrm{V}\) .)

\begin{enumerate}
\def\labelenumi{\alph{enumi})}
\setcounter{enumi}{2}
\tightlist
\item
  Prove the formula \(\int\overline{\mathrm{P}^{*}(ix)}\mathrm{Q}^{*}(ix)d\gamma(x)=(\mathrm{P}|\mathrm{Q})\) for P and Q in \(\mathbf{S}(\mathrm{V}^{*})\) (identified with a subspace of \(\mathbf{S}_{\mathbf{C}}((\mathrm{V}\otimes\mathbf{C})^{*})\) ).
\end{enumerate}

\(d)\) If \(\mathrm{P}\) is a homogeneous polynomial on \(\mathrm{V}\) such that \(\Delta (\mathrm{P}) = 0\) , then \(\int_{\mathrm{V}}\mathrm{P}(x)^2 d\gamma (x) = (\mathrm{P}|\mathrm{P})\) .

\begin{enumerate}
\def\labelenumi{\alph{enumi})}
\setcounter{enumi}{4}
\tightlist
\item
  Let W be a finite group of automorphisms of V generated by reflections; denote by H the set of reflections in W. For \(h \in H\) , let \(e_{h} \in V\) and \(f_{h} \in V^{*}\) be such that \(h(x) = x + f_{h}(x)e_{h}\) . Show that the polynomial \(P = \prod_{h \in H} f_{h}\) satisfies \(\Delta(P) = 0\) (use Chap. V, §5, no. 4, Prop. 5).
\end{enumerate}

\begin{enumerate}
\def\labelenumi{\arabic{enumi})}
\setcounter{enumi}{2}
\tightlist
\item
  Let \(\mu\) be a Haar measure on the additive group \(\mathfrak{g}\) .
\end{enumerate}

\begin{enumerate}
\def\labelenumi{\alph{enumi})}
\item
  There exists a unique Haar measure \(\mu_{G}\) on G with the following property: if \(\omega_{G}\) and \(\omega_{g}\) are invariant differential forms of degree n on G and g, respectively, such that \(\omega_{\mathrm{G}}(e)=\omega_{\mathfrak{g}}(0)\) and \(|\omega_{g}|=\mu\) , then \(|\omega_{G}|=\mu_{G}\) . The map \(\mu\mapsto\mu_{G}\) is a bijection from the set of Haar measures on g to the analogous set for G.
\item
  Choose a scalar product ( \textbar{} ) invariant under g, and a Haar measure \(\tau\) on t such that \(\mu\) (resp. \(\tau\) ) corresponds to the Lebesgue measure when g is identified with \(R^{n}\) (resp. t with \(R^{r}\) ) by means of an orthonormal basis. Prove the formula
\end{enumerate}

\[
\int_ {t} \pi_ {\mathfrak {g}} (x) e ^ {- (x | x) / 2} d \tau (x) = (2 \pi) ^ {n / 2} \frac {\tau_ {\mathrm{T}} (\mathrm{T})}{\mu_ {\mathrm{G}} (\mathrm{G})} w (\mathrm{G}).
\]

\begin{enumerate}
\def\labelenumi{\alph{enumi})}
\setcounter{enumi}{2}
\tightlist
\item
  Identify X(T) with a subset of the Hilbert space \(t^{*}\) by means of the map \((2\pi i)^{-1}\delta\) (§4, no. 2), and put \(\mathrm{P}(x)=\prod_{\alpha\in\mathrm{R}_{+}}\langle\alpha,x\rangle\) for \(x\in t\) . With the notations of Exerc. 1 and 2, show that
\end{enumerate}

\[
\frac {\tau_ {\mathrm{T}} (\mathrm{T})}{\mu_ {\mathrm{G}} (\mathrm{G})} = (2 \pi) ^ {\mathrm{N}} w (\mathrm{G}) ^ {- 1} (\mathrm{P} | \mathrm{P}) = \pi^ {\mathrm{N}} \prod_ {\alpha \in \mathrm{R} _ {+}} (\alpha | \alpha) \prod_ {i = 1} ^ {l} m _ {i}!
\]

\(d\) ) With the notations of Exerc. 9 of §3, let \(\mathfrak{g}_{\mathbf{Z}}\) be the Z-Lie subalgebra of \(\mathfrak{g}\) generated by \((2\pi)^{-1}\Gamma (\mathrm{T})\) and the elements \(u_{\alpha},v_{\alpha}\) for \(\alpha \in \mathbb{R}\) . Prove the equality

\[
\mu_ {\mathrm{G}} (\mathrm{G}) = \mu (\mathfrak {g} / \mathfrak {g} _ {\mathbf {Z}}) \frac {2 ^ {r} \pi^ {\mathrm{N} + r}}{\prod_ {i} m _ {i} !}.
\]

\begin{enumerate}
\def\labelenumi{\alph{enumi})}
\setcounter{enumi}{4}
\tightlist
\item
  Assume that G is simply-connected. Show that l = r and that
\end{enumerate}

\[
\mu_ {\mathrm{G}} (\mathrm{G}) = 2 ^ {l / 2} \pi^ {- \mathrm{N}} f ^ {1 / 2} \prod_ {\alpha \in \mathrm{R} _ {+}} (\alpha | \alpha) ^ {- 1} \prod_ {i = 1} ^ {l} (\alpha_ {i} | \alpha_ {i}) ^ {- 1 / 2} \prod_ {i} (m _ {i}!) ^ {- 1},
\]

where \(\{\alpha_1, \ldots, \alpha_l\}\) is a basis of \(R\) and \(f\) is the connection index of \(R\) .

\(f\) ) Assume further that \(\mathbb{R}\) is irreducible and that all its roots are of the same length; take the scalar product on \(\mathfrak{g}\) to be the opposite of the Killing form. Prove that \(\mu_{\mathrm{G}}(\mathrm{G}) = (2\pi)^{\mathrm{N} + r}(2h)^{n / 2}f^{1 / 2}\prod_{i}(m_i!)^{-1}\) .

\begin{enumerate}
\def\labelenumi{\arabic{enumi})}
\setcounter{enumi}{3}
\tightlist
\item
  Let \(\mathrm{X}\) be a differentiable manifold of class \(\mathrm{C}^{\infty}\) . In this and the following exercises, we shall denote simply by \(\mathrm{H}(\mathrm{X})\) the graded \(\mathbf{R}\) -space \(\mathrm{H}(\Omega (\mathrm{X}))\) .
\end{enumerate}

\begin{enumerate}
\def\labelenumi{\alph{enumi})}
\item
  Show that \(\Omega(\mathrm{X})\) is an associative and anti-commutative graded differential algebra (Algèbre, Chap. X, p.~183, Exerc. 18). Deduce that \(\mathrm{H}(\mathrm{X})\) has a natural associative and anti-commutative graded algebra structure.
\item
  If X is connected and of dimension p, then \(\mathrm{H}^{i}(\mathrm{X}) = 0\) for i \textgreater{} p and \(\dim_{\mathbf{R}} \mathrm{H}^{0}(\mathrm{X}) = 1\) . Moreover, if X is compact, the space \(\mathrm{H}^{p}(\mathrm{X})\) is of dimension one (use Differentiable and Analytic Manifolds, Results, 11.2.4).
\item
  Let Y be another manifold of class \(C^{\infty}\) , and \(f: X \to Y\) a map of class \(C^{\infty}\) . The map \(f^{*}: \Omega(Y) \to \Omega(X)\) is a morphism of complexes (Differentiable and Analytic Manifolds, Results, 8.3.5); show that \(\mathrm{H}(f^{*}) : \mathrm{H}(Y) \to \mathrm{H}(X)\) is a morphism of algebras.
\item
  Assume that we are given a law of operation of G on X of class \(C^{\infty}\) . Show that the subcomplex \(\Omega(\mathrm{X})^{\mathrm{G}}\) of invariant forms is a subalgebra of \(\Omega(\mathrm{X})\) . Deduce that the map \(\mathrm{H}(i):\mathrm{H}(\Omega(\mathrm{X})^{\mathrm{G}})\to\mathrm{H}(\mathrm{X})\) defined in Th. 2 is an isomorphism of graded algebras.
\item
  Show that \((\operatorname{Alt}(\mathfrak{g}))^{\mathrm{G}}\) is a subalgebra of \(\operatorname{Alt}(\mathfrak{g})\) , and that the graded algebra \(\operatorname{H}(\operatorname{G})\) is isomorphic to it.
\end{enumerate}

\begin{enumerate}
\def\labelenumi{\arabic{enumi})}
\setcounter{enumi}{4}
\tightlist
\item
  Denote by H(G) the graded R-algebra H(Ω(G)) (cf.~Exerc. 4).
\end{enumerate}

\(a\) ) The space \(\mathrm{H}^p (\mathbf{G})\) is zero for \(p > n\) and of dimension one for \(p = n\) and \(p = 0\) . The space \(\mathrm{H}^1 (\mathrm{G})\) can be identified canonically with \(\mathfrak{c}^*\) , where \(\mathfrak{c} = \mathrm{L}(\mathrm{C}(\mathrm{G}))\) .

\begin{enumerate}
\def\labelenumi{\alph{enumi})}
\setcounter{enumi}{1}
\item
  Assume from now on that G is semi-simple. Show that \(\mathrm{H}^{1}(\mathrm{G})=\mathrm{H}^{2}(\mathrm{G})=0\) (cf.~Chap. I, §6, Exerc. 1).
\item
  Denote by \(\mathrm{B}(\mathfrak{g})\) the space of G-invariant symmetric bilinear forms on \(\mathfrak{g}\) ; for \(b \in \mathrm{B}(\mathfrak{g})\) and \(x, y, z\) in \(\mathfrak{g}\) , put \(\tilde{b}(x, y, z) = b([x, y], z)\) . Show that the map \(b \mapsto \tilde{b}\) defines an isomorphism from \(\mathrm{B}(\mathfrak{g})\) to \(\mathrm{H}^{3}(\mathrm{G})\) (let \(\omega\in\mathrm{H}^{3}(\mathrm{G})\) ; prove that, for all \(x\in g\) , there exists a unique linear form \(f(x)\) on g such that \(df(x)=i(x)\omega\) , and consider the form \((x,y)\mapsto-\langle y,f(x)\rangle\) ).
\end{enumerate}

\(d\) ) Show that the dimension of the \(\mathbf{R}\) -vector space \(\mathrm{H}^3 (\mathrm{G})\) is equal to the number of simple ideals of \(\mathfrak{g}\) .

\begin{enumerate}
\def\labelenumi{\arabic{enumi})}
\setcounter{enumi}{5}
\tightlist
\item
  Put \(b_{i}(\mathrm{G}) = \dim_{\mathbf{R}} \mathrm{H}^{i}(\mathrm{G})\) for \(i \geq 0\) and, if \(\mathrm{X}\) denotes an indeterminate, \(\mathrm{P}_{\mathrm{G}}(\mathrm{X}) = \sum_{i \geq 0} b_{i}(\mathrm{G}) \mathrm{X}^{i}\) .
\end{enumerate}

\begin{enumerate}
\def\labelenumi{\alph{enumi})}
\item
  Prove that \(b_{i}(\mathrm{G}) = \int_{\mathrm{G}} \mathrm{Tr} \wedge^{i} (\mathrm{Ad} g) dg\) and \(\mathrm{P}_{\mathrm{G}}(\mathrm{X}) = \int_{\mathrm{G}} \det(1 + \mathrm{X}. \mathrm{Ad} g) dg\) (cf.~Appendix II, Lemma 1).
\item
  Deduce the equalities \(\sum_{i} b_{i}(\mathrm{G}) = 2^{r}\) , and \(\sum_{i} (-1)^{i} b_{i}(\mathrm{G}) = 0\) if \(\dim \mathrm{G} > 0\) (use H. Weyl's formula, or Exerc. 8 of §2).
\item
  Take \(\mathrm{G} = \mathbf{U}(n, \mathbf{C})\) . Show that \(\mathrm{P}_{\mathrm{G}}(\mathrm{X})\) is the coefficient of \((\mathrm{X}_1 \ldots \mathrm{X}_n)^{2n-2}\) in the polynomial \(\frac{1}{n!}(1 + \mathrm{X})^n \prod_{\substack{1 \leq i, j \leq n \\ i \neq j}} (\mathrm{XX}_i + \mathrm{X}_j)(\mathrm{X}_i - \mathrm{X}_j)\) (with coefficients in \(\mathbf{Z}[\mathrm{X}]\) ).
\end{enumerate}

\begin{enumerate}
\def\labelenumi{\arabic{enumi})}
\setcounter{enumi}{6}
\tightlist
\item
  Let \(\mathbf{K}\) be a connected compact Lie group.
\end{enumerate}

\begin{enumerate}
\def\labelenumi{\alph{enumi})}
\item
  Let \(f: K \to G\) be a surjective homomorphism with finite kernel. Show that the homomorphism \(\mathrm{H}(f^{*}) : \mathrm{H}(\mathrm{G}) \to \mathrm{H}(\mathrm{K})\) is an isomorphism.
\item
  Show that the algebra \(\mathrm{H}(\mathrm{G}\times\mathrm{K})\) can be identified canonically with the skew tensor product \(\mathrm{H}(\mathrm{G})^{g}\otimes\mathrm{H}(\mathrm{K})\) .
\item
  Deduce from a) and b) that \(\mathrm{H}(\mathrm{G})\) is isomorphic to \(\mathrm{H}(\mathrm{C}(\mathrm{G})_{0})^{g}\otimes\mathrm{H}(\mathrm{D}(\mathrm{G}))\) as an algebra. Show that the algebra \(\mathrm{H}(\mathrm{C}(\mathrm{G})_{0})\) is isomorphic to \(\wedge(\mathfrak{c}^{*})\) , with \(\mathfrak{c}=\mathrm{L}(\mathrm{C}(\mathrm{G}))\) .
\end{enumerate}

¶ 8) Let k be a field of characteristic zero and E a graded left bigebra over k (Algebra, Chap. III, §11, no. 4, Def. 3). Assume that E is anti-commutative and anti-cocommutative, and that \(E_{m} = 0\) for m sufficiently large. Denote by P the set of primitive elements of E (cf.~Chap. II, §1).

\begin{enumerate}
\def\labelenumi{\alph{enumi})}
\item
  Show that every homogeneous element of P is of odd degree (expand \(c(x^{m})\) for m large). Deduce that there is a canonical morphism \(\varphi : \Lambda(\mathrm{P}) \to \mathrm{E}\) of graded left bigebras (the graded left bigebra structure of \(\Lambda(\mathrm{P})\) being that defined in Algebra, Chap. III, §11, Exerc. 6).
\item
  Show that \(\varphi\) is an isomorphism (adapt the proof of Th. 1 of Chap. II, §1, no. 6).
\end{enumerate}

\begin{enumerate}
\def\labelenumi{\arabic{enumi})}
\setcounter{enumi}{8}
\tightlist
\item
  Denote by \(m: \mathrm{G} \times \mathrm{G} \to \mathrm{G}\) the map such that \(m(g, h) = gh\) for \(g, h\) in \(\mathrm{G}\) . Identify \(\mathrm{H}(\mathrm{G} \times \mathrm{G})\) with \(\mathrm{H}(\mathrm{G})^g \otimes \mathrm{H}(\mathrm{G})\) (Exerc. 7 b)), so that \(m^*\) defines a homomorphism of algebras \(c: \mathrm{H}(\mathrm{G}) \to \mathrm{H}(\mathrm{G})^g \otimes \mathrm{H}(\mathrm{G})\) .
\end{enumerate}

\begin{enumerate}
\def\labelenumi{\alph{enumi})}
\item
  Show that \((\mathrm{H}(\mathrm{G}), c)\) is an anti-commutative and anti-cocommutative graded left bigebra (observe that the map \(g \mapsto g^{-1}\) induces on \(\mathrm{H}^{p}(\mathrm{G})\) the multiplication by \((-1)^{p}\) ).
\item
  Let \(\mathrm{P}(\mathrm{G})\) be the graded subspace of \(\mathrm{H}(\mathrm{G})\) consisting of the primitive elements; deduce from Exerc. 8 that there is an isomorphism of graded bigebras \(\Lambda(\mathrm{P}(\mathrm{G})) \to \mathrm{H}(\mathrm{G})\) .
\item
  Show that \(\dim_{\mathbf{R}}\mathrm{P}(\mathrm{G})=r\) (use Exerc. 6 b)).
\end{enumerate}

\(d\) ) Deduce that the polynomial \(\sum_{i\geq 0}b_i(\mathrm{G})\mathrm{X}^i\) is of the form

\[
(1 + \mathrm{X}) ^ {c} \prod_ {i = 1} ^ {l} (1 + \mathrm{X} ^ {2 k _ {i} + 1}),
\]

where c is the dimension of \(\mathrm{C}(\mathrm{G})\) , l the rank of \(\mathrm{D}(\mathrm{G})\) , and the \(k_{i}\) are integers \(\geq 1\) ; we have \(k_{1} + \cdots + k_{l} = \frac{1}{2}\mathrm{Card}\mathrm{R}(\mathrm{G},\mathrm{T})\)\(^{12}\) .

\footnotetext[12]{The integers \(k_{i}\) are in fact the exponents of \(\mathrm{R}(\mathrm{G},\mathrm{T})\). Cf. J. LERAY, Sur l'homologie des groupes de Lie, des espaces homogènes et des espaces fibrés principaux, Colloque de Topologie de Bruxelles (1950), pp. 101-115.}

\begin{enumerate}
\def\labelenumi{\arabic{enumi})}
\setcounter{enumi}{9}
\tightlist
\item
  Let \(\mathrm{H}\) be a closed subgroup of \(\mathrm{G}\) ; denote by \(dh\) the Haar measure on \(\mathrm{H}\) of total mass 1. Put \(\chi(\mathrm{G} / \mathrm{H}) = \sum_{i > 0} (-1)^i \dim_{\mathbf{R}} \mathrm{H}^i (\mathrm{G} / \mathrm{H})\) .
\end{enumerate}

\begin{enumerate}
\def\labelenumi{\alph{enumi})}
\item
  Prove the equality \(\chi(\mathrm{G}/\mathrm{H})=\overline{\int_{\mathrm{H}}}\det(1-\mathrm{Ad}_{\mathfrak{g}/\mathfrak{h}}h)\,dh\) (use Algèbre, Chap. X, p.~41, Prop. 11 and Lemma 1 of Appendix II, no. 2).
\item
  Let \(\pi_{0}(\mathrm{H})\) be the number of connected components of H. Show that
\end{enumerate}

\(\chi (\mathrm{G / H}) = 0\) if \(\mathrm{H_0}\) is not of maximum rank

\(\chi (\mathrm{G / H}) = w(\mathrm{G}) / w(\mathrm{H}_0)\pi_0(\mathrm{H})\) if \(\mathrm{H_0}\) is of maximum rank.

\begin{enumerate}
\def\labelenumi{\arabic{enumi})}
\setcounter{enumi}{10}
\tightlist
\item
  Let \(u\) be an automorphism of \(G\) of order two; denote by \(K\) the identity component of the set of fixed points of \(u\) , \(\mathfrak{k}\) its Lie algebra, and \(X\) the symmetric space \(G / K\) (§1, Exerc. 8).
\end{enumerate}

\begin{enumerate}
\def\labelenumi{\alph{enumi})}
\item
  Show that every G-invariant differential form \(\omega\) on X satisfies \(d\omega = 0\) (observe that \(u\) induces on \(\mathrm{Alt}^p (\mathfrak{g} / \mathfrak{k})\) the multiplication by \((-1)^p)\) .
\item
  Deduce that there is an isomorphism from the graded algebra \(\mathrm{H}(\mathrm{G}/\mathrm{K})\) to the graded subalgebra of \(\operatorname{Alt}(\mathfrak{g}/\mathfrak{k})\) consisting of the K-invariant elements.
\item
  Put \(b_{i}(\mathrm{G / K}) = \dim_{\mathbf{R}}\mathrm{H}^{i}(\mathrm{G / K})\) for \(i\geq 0\) ; for \(k\in \mathrm{K}\) , denote by \(\mathrm{Ad}^{-}k\) the restriction of \(\mathrm{Ad}k\) to the eigenspace of \(\mathrm{L}(u)\) with eigenvalue \(-1\) . Let \(dk\) be the Haar measure on \(\mathbf{K}\) of total mass 1. Prove the formulas \(b_{i}(\mathrm{G / K}) = \int_{\mathrm{K}}\mathrm{Tr}\wedge^{i}(\mathrm{Ad}^{-}k)dk\) and \(\sum_{i > 0}b_{i}(\mathrm{G / K})\mathrm{X}^{i} = \int_{\mathrm{K}}\det (1 + \mathrm{XAd}^{-}k)dk\) .
\item
  Prove that if, in addition, K is of maximum rank, the algebra \(\mathrm{H}(\mathrm{G}/\mathrm{K})\) is zero in odd degrees (observe that \(u = \operatorname{Int} k\) , with \(k \in K\) ).
\item
  Calculate the graded algebra \(\mathrm{H}(\mathbf{S}_{n})\) . Deduce that \(S_{n}\) admits a Lie group structure (compatible with its manifold structure) if and only if n is equal to 1 or 3.
\end{enumerate}

\begin{enumerate}
\def\labelenumi{\arabic{enumi})}
\setcounter{enumi}{11}
\tightlist
\item
  Let H be a connected closed subgroup of G and h its Lie algebra.
\end{enumerate}

\begin{enumerate}
\def\labelenumi{\alph{enumi})}
\item
  Let \(\alpha\) be an element of \(\mathfrak{h}^*\) invariant under H; show that there exists an element \(\bar{\alpha}\) of \(\mathfrak{g}^*\) invariant under H whose restriction to \(\mathfrak{h}\) is equal to \(\alpha\) . Show that the element \(d\bar{\alpha} \in \mathrm{Alt}^2(\mathfrak{g})\) is annihilated by \(i(\eta)\) and \(\theta(\eta)\) for all \(\eta \in \mathfrak{h}\) , and that its class in \(\mathrm{H}^2(\mathrm{G}/\mathrm{H})\) does not depend on the choice of \(\bar{\alpha}\) . This defines a homomorphism \(\varphi: \mathrm{H}^1(\mathrm{H}) \to \mathrm{H}^2(\mathrm{G}/\mathrm{H})\) .
\item
  Assume from now on that \(\mathbf{G}\) is semi-simple. Prove that \(\varphi\) is an isomorphism.
\item
  Prove that \(\mathrm{H}\) is semi-simple if and only if \(\mathrm{H}^2 (\mathrm{G / H}) = 0\) .
\end{enumerate}

\(d\) ) Without assuming that \(\mathrm{H}\) is connected, define an isomorphism \((\mathfrak{h}^{*})^{\mathrm{H}}\to\) \(\mathrm{H}^2 (\mathrm{G / H})\) (apply \(b\) ) to \(\mathrm{H}_0\) ).

\begin{enumerate}
\def\labelenumi{\arabic{enumi})}
\setcounter{enumi}{12}
\tightlist
\item
  Let \(\mathrm{H}\) be a closed subgroup of \(\mathrm{G}\) ; put \(\mathrm{X} = \mathrm{G} / \mathrm{H}\) and \(n = \dim \mathrm{X}\) . Show that the following conditions are equivalent:
\end{enumerate}

\begin{enumerate}
\def\labelenumi{(\roman{enumi})}
\item
  There exists a 2-form \(\omega\) on X such that \(d\omega = 0\) and such that the alternating form \(\omega_{x}\) on \(\mathrm{T}_x(\mathrm{X})\) is separating for all \(x\in \mathrm{X}\) ;
\item
  \(n\) is even, and there exists an element \(\omega\) of \(\mathrm{H}^2 (\mathrm{G / H})\) such that \(\omega^{n / 2}\neq 0\) ; (iii) \(\mathrm{H}\) is the centralizer of a torus in \(\mathbf{G}\) ;
\item
  H is of maximum rank, and there exists a G-invariant complex structure on X;
\item
  there exists a complex structure \(j\) and a 2-form \(\omega\) on X such that \(d\omega = 0\) , \(\omega_{x}(ju,jv) = \omega_{x}(u,v)\) and \(\omega_{x}(u,ju) > 0\) for all \(x\) in X and all non-zero \(u,v\) in \(\mathrm{T}_x(\mathrm{X})\) (* in other words, a Kähler structure on \(\mathrm{X}_{*}\) );
\item
  there exists a complex structure j and a 2-form \(\omega\) on X satisfying the conditions in (v) and invariant under G (* that is, a G-invariant Kähler structure on \(X_{*}\) ).
\end{enumerate}

(To prove that (ii) \(\Longrightarrow\) (iii), put \(S = C(H)_{0}\) and \(Z = Z_{G}(S)\) ; by using Exerc. 12 d), show that the canonical map \(H^{2}(G/Z) \to H^{2}(G/H)\) is surjective, and deduce that Z = H. The equivalence of (iii) and (iv) follows from Exerc. 8 of §4. To prove that (iii) \(\Longrightarrow\) (vi), construct an H-invariant separating positive hermitian form on g/L(H) and consider its imaginary part; use Exerc. 11 a).

\exercisesection{\S~7}

\begin{enumerate}
\def\labelenumi{\arabic{enumi})}
\tightlist
\item
  Take G to be the group \(\mathbf{U}(n,\mathbf{C})\) , and T to be the subgroup consisting of the diagonal matrices; for \(t = \text{diag}(t_1, \ldots, t_n)\) in T and \(1 \leq i \leq n\) , put \(\varepsilon_i(t) = t_i\) . Denote by \(\sigma\) the identity representation of G on \(C^n\) .
\end{enumerate}

\begin{enumerate}
\def\labelenumi{\alph{enumi})}
\item
  The group X(T) has basis \(\varepsilon_{1},\ldots,\varepsilon_{n}\) ; show that every element of \(\mathbf{Z}[\mathrm{X}(\mathrm{T})]^{\mathrm{W}}\) is of the form \(e^{k(\varepsilon_{1}+\cdots+\varepsilon_{n})}\mathrm{P}(e^{\varepsilon_{1}},\ldots,e^{\varepsilon_{n}})\) , where k is a relative integer and P a symmetric polynomial in n variables with integer coefficients.
\item
  Show that the representations \(\Lambda^r\sigma\) ( \(1 \leq r \leq n\) ) are irreducible.
\item
  Show that the homomorphism \(u: \mathbf{Z}[X_1, X_2, \ldots, X_n][X_n^{-1}] \to \mathrm{R}(\mathrm{G})\) such that \(u(X_i) = [\Lambda^i \sigma]\) is an isomorphism.
\end{enumerate}

\begin{enumerate}
\def\labelenumi{\arabic{enumi})}
\setcounter{enumi}{1}
\tightlist
\item
  Let \(\mathrm{G} = \mathbf{SO}(2l + 1, \mathbf{R})\) , with \(l \geq 1\) ; we use the notations of Exerc. 2 of §4. Denote by \(\sigma\) the representation of \(\mathrm{G}\) on \(\mathbf{C}^{2l + 1}\) obtained from the identity representation by extension of scalars.
\end{enumerate}

\begin{enumerate}
\def\labelenumi{\alph{enumi})}
\item
  The Z-module X(T) has basis \(\varpi_{1},\ldots,\varpi_{l-1},2\varpi_{l}\) , as well as \(\varepsilon_{1},\ldots,\varepsilon_{l}\) .
\item
  Denote by \(\eta_{i}\) the element \(e^{\varepsilon_{i}} + e^{-\varepsilon_{i}}\) of Z{[}X(T){]}. Show that every element of \(\mathbf{Z}[X(T)]^{W}\) can be written as \(\mathrm{P}(\eta_{1}, \ldots, \eta_{l})\) where P is a symmetric polynomial in l variables with integer coefficients.
\item
  Show that the representations \(\Lambda^r\sigma (r\leq 2l + 1)\) are irreducible. For \(r\leq l\) , prove the equality
\end{enumerate}

\[
\begin{array}{c} \operatorname{Ch} (\Lambda^ {r}   \sigma) = s _ {r} (\eta_ {1}, \ldots , \eta_ {l}) + (l - r + 2) s _ {r - 2} (\eta_ {1}, \ldots , \eta_ {l}) + \dots + \\ \qquad \qquad + \binom {l - r + 2 k} {k} s _ {r - 2 k} (\eta_ {1}, \ldots , \eta_ {l}) + \dots , \end{array}
\]

where the \(s_{k}\) are the elementary symmetric polynomials in l variables.

\(d\) ) Show that the homomorphism \(u:\mathbf{Z}[\mathrm{X}_1,\dots ,\mathrm{X}_l]\to \mathrm{R}(\mathrm{G})\) such that \(u(\mathrm{X}_i) = [\Lambda^i\sigma ]\) is an isomorphism.

\begin{enumerate}
\def\labelenumi{\arabic{enumi})}
\setcounter{enumi}{2}
\tightlist
\item
  Let \(\mathrm{G} = \mathbf{SO}(2l, \mathbf{R})\) , with \(l \geq 2\) ; we use the notations of Exerc. 2 of §4. Denote by \(\sigma\) the representation of G on \(C^{2l}\) obtained from the identity representation by extension of scalars.
\end{enumerate}

\begin{enumerate}
\def\labelenumi{\alph{enumi})}
\item
  The \(\mathbf{Z}\) -module \(\mathrm{X(T)}\) has basis \(\varepsilon_1, \ldots, \varepsilon_l\) .
\item
  Denote by \(\eta_i\) the element \(e^{\varepsilon_i} + e^{-\varepsilon_i}\) of \(\mathbf{Z}[\mathrm{X(T)}]\) and by \(\delta\) the element \(\prod_{i=1}^{l}(e^{\varepsilon_i} - e^{-\varepsilon_i})\) . Show that every element of \(\mathbf{Z}[\mathrm{X(T)}]^{\mathrm{W}}\) can be written as \(\mathrm{P}(\eta_1, \ldots, \eta_l) + \mathrm{Q}(\eta_1, \ldots, \eta_l)\delta\) , where \(\mathrm{P}\) and \(\mathrm{Q}\) are symmetric polynomials in \(l\) variables with coefficients in \(\frac{1}{2}\mathbf{Z}\) .
\item
  Show that the representations \(\Lambda^{r}\sigma\) are irreducible for \(r\neq l\) ; the representation \(\Lambda^{l}\sigma\) is the direct sum of two subrepresentations \(\tau^{+}\) and \(\tau^{-}\) , of highest weights \(2\varpi_{l}\) and \(2\varpi_{l-1}\) , respectively (see Chap. VIII, §13, Exerc. 10).
\end{enumerate}

\(d\) ) Show that, for \(r \leq l\) , the element \(\operatorname{Ch}(\Lambda^r \sigma)\) is given by the formula in Exerc. 2 \(c\) ) and that

\[
\mathrm{Ch} (\tau^ {+}) = \frac {1}{2} (\delta + \mathrm{Ch} (\Lambda ^ {l} \sigma)) \qquad \mathrm{Ch} (\tau^ {-}) = \frac {1}{2} (- \delta + \mathrm{Ch} (\Lambda ^ {l} \sigma)).
\]

\begin{enumerate}
\def\labelenumi{\alph{enumi})}
\setcounter{enumi}{4}
\tightlist
\item
  Show that the homomorphism \(u : Z[X_{1}, \ldots, X_{l}; Y] \to R(G)\) such that \(u(X_{i}) = [\Lambda^{i} \sigma]\) and \(u(Y) = [\tau^{+}]\) is surjective, and that its kernel is generated by \(Y^{2} - YX_{l} + A\) , with \(A \in Z[X_{1}, \ldots, X_{l}]\) .
\end{enumerate}

\begin{enumerate}
\def\labelenumi{\arabic{enumi})}
\setcounter{enumi}{3}
\tightlist
\item
  Let E be a complex vector space, and let \(\mathbf{T}_{q}^{p}(\mathrm{E})\) be the space of tensors of type \((p,q)\) on E (Algebra, Chap. III, §5, no. 6). Denote by \(\mathrm{H}_{q}^{p}(\mathrm{E})\) the subspace of \(\mathbf{T}_{q}^{p}(\mathrm{E})\) consisting of the symmetric tensors (that is, those belonging to the image of \(\mathbf{T}\mathbf{S}^{p}(\mathrm{E})\otimes\mathbf{T}\mathbf{S}^{q}(\mathrm{E}^{*})\) in \(\mathrm{T}_{q}^{p}(\mathrm{E})\) and annihilated by the contractions \(c_{j}^{i}\) for \(i\in\{1,p\}\) and \(j\in\{p+1,p+q\}\) (loc. cit.).
\end{enumerate}

\begin{enumerate}
\def\labelenumi{\alph{enumi})}
\item
  Show that \(\mathrm{H}_{q}^{p}(\mathbf{C}^{n})\) is stable under \(\mathbf{GL}(n,\mathbf{C})\) , and consequently under \(\mathbf{SU}(n,\mathbf{C})\) ; denote this representation of \(\mathbf{SU}(n,\mathbf{C})\) by \(\tau_{q}^{p}\) .
\item
  Show that \(\tau_q^p\) is an irreducible representation whose highest weight (with the notations of Exerc. 1 of §4) is \(p\varpi_1 + q\varpi_{n - 1}\) .
\item
  Every irreducible representation of \(\mathbf{SU}(3,\mathbf{C})\) is isomorphic to one of the representations \(\tau_q^p\) .
\end{enumerate}

\(d\) ) Let \((x_{i})_{1\leq i\leq n}\) be the canonical basis of \(\mathbf{C}^n\) , \((y_{j})_{1\leq j\leq n}\) the dual basis. Show that \(\mathrm{H}_q^p (\mathbf{C}^n)\) can be identified with the space of polynomials \(\mathrm{P}\in \mathbf{C}[x_1,\ldots ,x_n,y_1,\ldots ,y_n]\) , homogeneous of degree \(p\) in the \(x_{i}\) and of degree \(q\) in the \(y_{i}\) , and such that \(\sum_{i = 1}^{n}\frac{\partial^2\mathrm{P}}{\partial x_i\partial y_i} = 0\) .

\begin{enumerate}
\def\labelenumi{\arabic{enumi})}
\setcounter{enumi}{4}
\tightlist
\item
  Let \(k\) be a commutative field of characteristic zero, \(V\) a vector space over \(k\) , and \(\varPsi\) a separating quadratic form on \(V\) . Let \(\varGamma \in \mathbf{S}^2(V^*)\) (resp. \(\varGamma^* \in \mathbf{S}^2(V)\) ) be the element associated to \(\varPsi\) (resp. the inverse form of \(\varPsi\) ). Denote by \(Q\) the endomorphism of \(S(V)\) given by taking the product with \(\varGamma^*\) , by \(\varDelta\) the endomorphism of \(S(V)\) given by taking the inner product with \(\varGamma\) , and by \(h\) the endomorphism of \(S(V)\) that reduces on \(S^r(V)\) to multiplication by \(-\frac{n}{2} - r\) . \(a\) ) If \((x_i)_{1\leq i\leq n}\) is an orthonormal basis of \(V\) , then \(Q(P) = \frac{1}{2}\left(\sum x_i^2\right).P\) and \(\varDelta(P) = \frac{1}{2}\sum \frac{\partial^2P}{\partial x_i^2}\) for \(P \in S(V)\) .
\end{enumerate}

\begin{enumerate}
\def\labelenumi{\alph{enumi})}
\setcounter{enumi}{1}
\item
  Prove the formulas \([\Delta, \mathrm{Q}] = -h\) , \([h, \Delta] = 2\Delta\) , \([h, \mathrm{Q}] = -2\mathrm{Q}\) .
\item
  Let \(H_{r}\) be the subspace of \(\mathbf{S}^{r}(\mathrm{V})\) annihilated by \(\Delta\) (``harmonic homogeneous polynomials of degree r''). Deduce from b) that there is a direct sum decomposition, stable under \(\mathbf{O}(\varPsi)\) , given by
\end{enumerate}

\[
\mathbf {S} ^ {r} (\mathrm{V}) = \mathrm{H} _ {r} \oplus \mathrm{QH} _ {r - 2} \oplus \mathrm{Q} ^ {2} \mathrm{H} _ {r - 4} \oplus \dots .
\]

\(d\) ) Show that the representation \(\mathrm{H}_r\) is irreducible (cf.~Chap. VIII, §13, no. 3, (IV)).

\begin{enumerate}
\def\labelenumi{\alph{enumi})}
\setcounter{enumi}{4}
\tightlist
\item
  Take \(k = \mathbf{C}\) , and \(V = \mathbf{C}^n\) equipped with the usual quadratic form ( \(n \geq 3\) ). We thus obtain irreducible representations \(\tau_r\) of \(\mathbf{SO}(n, \mathbf{R})\) ; with the notations of Exerc. 2 of §4, show that the highest weight of \(\tau_r\) is \(r\varpi_1\) .
\end{enumerate}

\(f\) ) Let \(\Gamma\) be the Casimir element of G obtained from the Killing form. Prove the formula \(\Gamma_{\mathbf{S}(\mathrm{V})} = \frac{1}{2n - 4}\left(-4\mathrm{Q}\Delta +\left(\mathrm{H} + \frac{n}{2}\mathrm{I}\right)\left(\mathrm{H} + \left(2 - \frac{n}{2}\right)\mathrm{I}\right)\right)\) . Calculate \(\tilde{\Gamma} (\tau_r)\) and deduce the value of the form \(Q_{\Gamma}\) (Prop. 4).

\begin{enumerate}
\def\labelenumi{\arabic{enumi})}
\setcounter{enumi}{5}
\item
  Assume that G is almost simple. Show that G admits a faithful irreducible representation if and only if it is not isomorphic to \(\mathbf{Spin}(4k, \mathbf{R})\) for \(k \geq 2\) .
\item
  Let \(\tau : \mathrm{G} \to \mathbf{SO}(n, \mathbf{R})\) be a real unitary representation of \(\mathrm{G}\) ; denote by \(\varphi : \mathbf{Spin}(n, \mathbf{R}) \to \mathbf{SO}(n, \mathbf{R})\) the canonical double covering. Then \(\tau\) is said to be spinorial if there exists a morphism \(\tilde{\tau} : \mathrm{G} \to \mathbf{Spin}(n, \mathbf{R})\) such that \(\varphi \circ \tilde{\tau} = \tau\) .
\end{enumerate}

\begin{enumerate}
\def\labelenumi{\alph{enumi})}
\item
  Let \(\Sigma\) be a subset of \(\mathrm{P}(\mathrm{T},\tau)\) such that \(\Sigma\cup(-\Sigma)=\mathrm{P}(\mathrm{T},\tau)\) and \(\Sigma\cap(-\Sigma)=\varnothing\) ; denote by \(\omega_{\Sigma}\) the sum of the elements of \(\Sigma\) . The class \(\varpi\) of \(\omega_{\Sigma}\) in \(\mathrm{X}(\mathrm{T})/2\mathrm{X}(\mathrm{T})\) is independent of the choice of \(\Sigma\) . Prove that \(\tau\) is spinorial if and only if \(\varpi=0\) .
\item
  Prove that \(\rho \in \mathrm{X}(\mathrm{T})\) if and only if the adjoint representation is spinorial.
\end{enumerate}

\begin{enumerate}
\def\labelenumi{\arabic{enumi})}
\setcounter{enumi}{7}
\tightlist
\item
  Let G be a Lie group whose Lie algebra is compact, and which has a finite number of connected components. Show that G has a faithful linear representation on a finite dimensional vector space (write G as the semi-direct product of a compact group K by a vector group N; choose a faithful representation of K on a finite dimensional vector space W and represent G as a subgroup of the affine group of \(W \oplus N\) ).
\end{enumerate}

\exercisesection{\S~8}

\begin{enumerate}
\def\labelenumi{\arabic{enumi})}
\tightlist
\item
  Let \(\mathbf{G} = \mathbf{SU}(2, \mathbf{C})\) , and let \(\sigma\) be the identity representation of G on \(C^{2}\) .
\end{enumerate}

\begin{enumerate}
\def\labelenumi{\alph{enumi})}
\item
  The irreducible representations of G are the representations \(\tau^{n} = S^{n}\sigma\) for \(n \geq 0\) .
\item
  Let \(e_{1}, e_{2}\) be the canonical basis of \(C^{2}\) ; show that the coefficients of \(\tau^{n}\) in the basis \((e_{1}^{i}e_{2}^{n-i})_{0\leq i\leq n}\) are the functions \(\tau_{ij}^{n}\) such that, for \(g=\begin{pmatrix}\alpha&\beta\\-\bar{\beta}&\bar{\alpha}\end{pmatrix}\in G\) , we have \(\tau_{ij}^{n}(g)=\frac{(-1)^{i}}{j!}\alpha^{i+j-n}\beta^{j-i}\mathrm{P}_{ij}^{n}(|\alpha|^{2})\) , with \(\mathrm{P}_{ij}^{n}(t)=\frac{d^{j}}{dt^{j}}\left[t^{n-i}(1-t)^{i}\right]\) (``Jacobi polynomials'').
\item
  Deduce from \(b)\) that the functions \((n + 1)^{1/2}\left(\frac{j!(n - j)!}{i!(n - i)!}\right)^{1/2}\tau_{ij}^n\) , for \(i,j,n\) integers with \(0\leq i\leq n\) , \(0\leq j\leq n\) , form an orthonormal basis of \(\mathrm{L}^2 (\mathrm{G})\) .
\end{enumerate}

\(d)\) For \(g = \left( \begin{array}{cc}\alpha & \beta \\ -\bar{\beta} & \bar{\alpha} \end{array} \right)\in \mathrm{G}\) , put \(\alpha = t^{1 / 2}e^{i(\varphi -\psi) / 2},\beta = (1 - t)^{1 / 2}e^{i(\varphi +\psi) / 2}\) , with \(0\leq t\leq 1,0\leq \varphi < 2\pi , - 2\pi \leq \psi < 2\pi\) . Show that the normalized Haar measure \(dg\) on \(\mathrm{G}\) is equal to \((8\pi^2)^{-1}dtd\varphi d\psi\) .

\begin{enumerate}
\def\labelenumi{\alph{enumi})}
\setcounter{enumi}{4}
\tightlist
\item
  Let \(a, b\) in \(\frac{1}{2}\mathbf{Z}\) be such that \(a - b \in \mathbf{Z}\) . Deduce from \(d\) ) that the polynomials \(\mathrm{P}_{n/2 - a, n/2 - b}^n(t)\) , for \(n\) an integer of the same parity as \(2a\) , \(n \geq \max(2a, 2b)\) , form an orthogonal basis of \(\mathrm{L}^2([0,1])\) for the measure \(t^{-a - b}(1 - t)^{a - b}dt\) .
\end{enumerate}

\begin{enumerate}
\def\labelenumi{\arabic{enumi})}
\setcounter{enumi}{1}
\item
  Let \(f\) be a complex-valued function on \(G\) of class \(C^\infty\) . Prove that there exist two complex-valued functions \(g\) and \(\varphi\) on \(G\) , of class \(C^\infty\) , such that \(\varphi\) is central and \(f = g * \varphi\) .
\item
  Let u be an irreducible representation of G, and let \(\lambda\) be its highest weight. For \(x \in g\) , denote by \(\bar{x}\) the unique element of \(\overline{C}\) that is conjugate to x under \(\mathrm{Ad}(G)\) . Prove the equality \(\|u(x)\|_{\infty} = |\delta(\lambda)(\bar{x})|\) .
\item
  Choose a separating positive quadratic form \(\mathbf{Q}\) on \(\mathfrak{g}\) invariant under \(\operatorname{Ad}(\mathbf{G})\) . For \(x \in \mathfrak{t}\) , put \(\vartheta_0(x) = \sum_{u \in \Gamma(\mathbf{T})} e^{-\mathbf{Q}(x + u)}\) .
\end{enumerate}

\begin{enumerate}
\def\labelenumi{\alph{enumi})}
\item
  Show that \(\vartheta_0\) is a function of class \(C^\infty\) on \(\mathfrak{t}\) , and that there exists a function \(\vartheta_1\) of class \(C^\infty\) on \(T\) such that \(\vartheta_1 \circ \exp_T = \vartheta_0\) .
\item
  Show that there exists a unique central function \(\vartheta\) on \(\mathbf{G}\) , of class \(\mathrm{C}^{\infty}\) , whose restriction to \(\mathrm{T}\) is equal to \(\vartheta_{1}\) . For every maximal torus \(\mathrm{S}\) of \(\mathrm{G}\) and all \(x\in \mathrm{L}(\mathrm{S})\) , we have
\end{enumerate}

\[
\vartheta (\exp x) = \sum_ {u \in \Gamma (S)} e ^ {- Q (x + u)}.
\]

\begin{enumerate}
\def\labelenumi{\alph{enumi})}
\setcounter{enumi}{2}
\tightlist
\item
  Let A be an alcove of t, dx a Haar measure on t, and h a locally integrable function on t invariant under the affine Weyl group \(W_{a}^{\prime}\) (§5, no. 2). Prove the equality
\end{enumerate}

\[
\int_ {\mathrm{A}} h (x) d x = \frac {1}{w (\mathrm{G})} \int_ {\mathfrak {t}} h (x) e ^ {- \mathrm{Q} (x)} (\vartheta_ {0} (x)) ^ {- 1} d x.
\]

\(d\) ) For \(x\in \mathfrak{g}\) , put \(\xi (x) = \lambda_{\mathfrak{g}}(x)e^{-\mathrm{Q}(x)}(\vartheta (\exp x))^{-1}\) , with \(\lambda_{\mathfrak{g}}(x) = \det \frac{e^{\operatorname*{ad}x} - 1}{\operatorname*{ad}x}\) (§6, no. 3). Show that \(\xi\) is a function of class \(\mathrm{C}^\infty\) on \(\mathfrak{g}\) and that if \(\mu\) is a Haar measure on \(\mathfrak{g}\) , the image under \(\exp_{\mathrm{G}}\) of the measure \(\xi \mu\) is a Haar measure on G (use \(c\) ) as well as Cor. 2 of Th. 1 and Prop. 4 (§6, no. 2 and 3)).

\begin{enumerate}
\def\labelenumi{\alph{enumi})}
\setcounter{enumi}{4}
\tightlist
\item
  Prove the formula \(\vartheta_0(x) = m\sum_{\lambda \in \mathrm{X(T)}}\exp (\delta (\lambda)(x) - \frac{1}{4}\mathbf{Q}'(\delta (\lambda)))\) for \(x\in \mathfrak{t}\) , where \(\mathbf{Q}'\) is the quadratic form on \(\mathfrak{t}_{\mathbf{C}}^{*}\) inverse to the quadratic form on \(\mathfrak{t}_{\mathbf{C}}\) induced by \(\mathbf{Q}\) and where \(m\) is a constant that should be calculated (use Poisson's formula, cf.~Spectral Theory, Chap. II, §1, no. 8). Deduce the equality \(\vartheta (t) = m\sum_{\lambda \in \mathrm{X(T)}}e^{-\mathrm{Q}'(\delta (\lambda)) / 4}t^{\lambda}\) for \(t\in \mathrm{T}\) .
\end{enumerate}

\begin{enumerate}
\def\labelenumi{\arabic{enumi})}
\setcounter{enumi}{4}
\tightlist
\item
  \begin{enumerate}
  \def\labelenumii{\alph{enumii})}
  \tightlist
  \item
    Let V be a real vector space, f a non-zero linear form on V, H the kernel of f.~Then, a function \(\varphi\) on V of class \(C^{\infty}\) vanishes on H if and only if it can be written as \(f\varphi'\) , where \(\varphi'\) is a function of class \(C^{\infty}\) on V.
  \end{enumerate}
\end{enumerate}

\begin{enumerate}
\def\labelenumi{\alph{enumi})}
\setcounter{enumi}{1}
\item
  Let V be a real vector space, \((f_{i})_{i\in\mathrm{I}}\) a finite family of non-zero linear forms such that the \(H_{i}=\operatorname{Ker}f_{i}\) are pairwise distinct. Then, a function \(\varphi\) on V of class \(C^{\infty}\) vanishes on the union of the \(H_{i}\) if and only if it can be written as \(\psi\prod_{i\in I}f_{i}\) , where \(\psi\) is a function on V of class \(C^{\infty}\) .
\item
  Let T be a torus, \((\alpha_{i})_{i\in I}\) a finite family of characters of T distinct from 1, such that the \(K_{i} = Ker \alpha_{i}\) are pairwise distinct. Then, a function \(\varphi\) on T of class \(C^{\infty}\) vanishes on the union of the \(K_{i}\) if and only if it can be written as \(\psi \cdot \prod_{i\in I}(\alpha_{i}-1)\) , where \(\psi\) is a function of class \(C^{\infty}\) on T (argue locally on T and use b)).
\end{enumerate}

\(d\) ) With the notations of no. 4, prove that the map \(b_{\infty}:\mathcal{C}^{\infty}(\mathrm{T})^{\mathrm{W}}\to \mathcal{C}^{\infty}(\mathrm{T})^{-\mathrm{W}}\) is bijective.

\begin{enumerate}
\def\labelenumi{\arabic{enumi})}
\setcounter{enumi}{5}
\item
  Assume that G is not commutative. Show that the continuous function \(\mathrm{J}(\rho)^{1/3}\) on T is anti-invariant under W, but does not belong to the image of the map \(b_{c}: \mathcal{C}(\mathrm{T})^{\mathrm{W}} \to \mathcal{C}(\mathrm{T})^{-\mathrm{W}}\) .
\end{enumerate}

\exercisesection{\S~9}

\begin{enumerate}
\def\labelenumi{\arabic{enumi})}
\item
  Let A be the compact subset of R consisting of 0 and the real numbers \(1/n\) , with \(n\) an integer \(\geq 1\) . Show that when \(\mathcal{C}^r(\mathbf{R};\mathbf{R})\) is given the topology of uniform \(\mathrm{C}^r\) -convergence on A, the set of morphisms that are embeddings in the neighbourhood of A is not open in \(\mathcal{C}^r(\mathbf{R};\mathbf{R})\) (consider a sequence of functions \((f_n)_{n\geq 1}\) such that \(f_n(x) = x\) for \(x\leq \frac{1}{n + 1}\) , \(f_n(x) = x - \frac{1}{n}\) for \(x\geq \frac{1}{n}\) ).\\
\item
  Let X be a separated manifold of class \(\mathrm{C}^r\) ( \(1\leq r\leq \infty\) ), countable at infinity, and pure of dimension \(n\) .
\end{enumerate}

\begin{enumerate}
\def\labelenumi{\alph{enumi})}
\item
  Assume that there exists an embedding \(\varphi\) of X into a finite dimensional vector space V. Show that there exists an embedding of X in \(\mathbf{R}^{2n + 1}\) (if \(\dim \mathrm{V} > 2n + 1\) , prove that there exists a point \(p\) of V such that, for all \(x \in \mathrm{X}\) , the straight line joining \(p\) to \(\varphi(x)\) meets \(\varphi(\mathrm{V})\) only in the point \(\varphi(x)\) , and meets it transversally at this point; deduce that there is an embedding of X in a space of dimension equal to \(\dim \mathrm{V} - 1\) ).
\item
  Show that there exists an embedding of X in \(R^{2n+1}\) . (Let O be the set of open subsets of X, U the subset of O consisting of the open sets U such that there exists a morphism \(\varphi: X \to R^{2n+1}\) whose restriction to U is an embedding; by using a), show that U is a quasi-full subset (General Topology, Chap. IX, §4, Exerc. 27) of O, and hence is equal to O.)
\item
  Show that there exists a proper embedding of X in \(\mathbf{R}^{2n + 1}\) . (Construct a proper embedding of X in \(\mathbf{R}^{2n + 2}\) by means of a proper function on X.)
\end{enumerate}

¶ 3) Let G be a Lie group, H a compact subgroup of G. Assume that G admits a faithful (finite dimensional) linear representation.

\begin{enumerate}
\def\labelenumi{\alph{enumi})}
\item
  Let \(\Theta_{\mathrm{H}}(\mathrm{G})\) be the subalgebra of \(\mathcal{C}(\mathrm{H};\mathbf{R})\) consisting of the restrictions to H of the (continuous) representative functions on G. Show that \(\Theta_{\mathrm{H}}(\mathrm{G})\) is dense in \(\mathcal{C}(\mathrm{H};\mathbf{R})\) for the topology of uniform convergence.
\item
  Let \(f \in \Theta_{\mathrm{H}}(\mathrm{G})\) . Show that there exists a representation \(\sigma\) of G on a finite dimensional real vector space such that a coefficient of the representation \(\sigma|H\) is not orthogonal to f (Spectral Theory, in preparation).
\item
  Let \(\rho:H\to\mathbf{GL}(V)\) be a representation of H on a finite dimensional real vector space. Show that there exists a finite dimensional real vector space W, a representation \(\sigma:G\to\mathbf{GL}(W)\) and an injective homomorphism \(u:V\to W\) , such that \(u(\rho(h)v)=\sigma(h)u(v)\) for \(h\in H\) , \(v\in V\) .
\end{enumerate}

(Reduce to the case in which \(\rho\) is irreducible, and use b).)

\begin{enumerate}
\def\labelenumi{\arabic{enumi})}
\setcounter{enumi}{3}
\item
  Let G be a Lie group, H a compact subgroup of G, \(\rho\) a unitary representation of H on a real Hilbert space V. Prove that there exists a real Hilbert space W, a unitary representation \(\sigma\) of G on W and an isometric injection \(u: V \to W\) with closed image, such that \(u(\rho(h)v) = \sigma(h)u(v)\) for \(h \in H\) , \(v \in V\) (same method as for Exerc. 3).
\item
  Let G be a Lie group, H a compact subgroup of G. Assume that there exists a faithful linear representation \(\rho: G \to \mathbf{GL}(V)\) of G on a finite dimensional real vector space.
\end{enumerate}

\begin{enumerate}
\def\labelenumi{\alph{enumi})}
\item
  Show that there exists a representation \(\sigma\) of \(\mathbf{GL}(\mathrm{V})\) on a finite dimensional real vector space W, a separating positive quadratic form q on V and a vector w in W such that \(\rho(\mathrm{H}) = \mathbf{O}(q) \cap \mathrm{F}_{w}\) , where \(F_{w}\) is the fixer of w in \(\mathbf{GL}(\mathrm{V})\) (choose a quadratic form q invariant under H and a representation of \(\mathbf{O}(q)\) such that \(\rho(\mathrm{H})\) is the fixer of a point (Cor. 2 of no. 2), then apply Exerc. 3c)).
\item
  Deduce from a) that there exists a finite dimensional real vector space E, a representation of G on E and a vector \(e \in E\) with fixer H (take \(E = W \oplus Q\) , where Q is the space of quadratic forms on V, and \(e = (w, q)\) ).
\end{enumerate}

\begin{enumerate}
\def\labelenumi{\arabic{enumi})}
\setcounter{enumi}{5}
\item
  Let G be a Lie group, H a compact subgroup of G. Show that there exists a unitary representation of G on a real Hilbert space E and a vector \(e \in E\) with fixer H (take \(\mathrm{E} = \mathrm{L}^{2}(\mathrm{G})\) ).
\item
  Let G be a Lie group, H a compact subgroup of G, and \(\rho\) a unitary representation of H on a real Hilbert space W. Denote by X the manifold \(G \times ^{H}W\) (no. 3).
\end{enumerate}

\begin{enumerate}
\def\labelenumi{\alph{enumi})}
\item
  Show that there exists a Hilbert space V, a unitary representation \(\sigma\) of G on V and an (analytic) embedding \(\varphi: X \to V\) such that \(\varphi(gx) = \sigma(g)\varphi(x)\) for \(g \in G\) , \(x \in X\) (use Exerc. 4 and Exerc. 6).
\item
  Prove that, if W is finite dimensional and if G admits a faithful finite dimensional linear representation, then V can be chosen to be finite dimensional (use Exerc. 3 and Exerc. 5).
\end{enumerate}

¶ 8) Let X be a paracompact manifold of class \(C^{r}\) ( \(1 \leq r \leq \infty\) ), and G a Lie group operating properly on X such that the law of operation \((g, x) \mapsto gx\) is of class \(C^{r}\) .

\begin{enumerate}
\def\labelenumi{\alph{enumi})}
\item
  Show that there exists a unitary representation \(\rho\) of G on a Hilbert space V and an embedding \(\varphi\) (of class \(C^{r}\) ) of X in V, such that \(\varphi(gx) = \sigma(g)\varphi(x)\) for \(g \in G, x \in X\) . (Use Prop. 6, Exerc. 7, and argue as in the proof of Prop. 4.)
\item
  Make the following additional assumptions:
\end{enumerate}

\begin{enumerate}
\def\labelenumi{(\roman{enumi})}
\item
  G admits a finite dimensional linear representation;
\item
  \(\mathrm{X}\) is countable at infinity, and of bounded dimension;
\item
  X has only a finite number of orbit types for the operation of G.
\end{enumerate}

Prove that there exists a finite covering of X by open sets \((\mathrm{U}_{i})_{i\in\mathrm{I}}\) that are stable under G, compact subgroups \((\mathrm{H}_{i})_{i\in\mathrm{I}}\) of G, and for all \(i\in I\) a closed submanifold \(S_{i}\) of \(U_{i}\) , stable under \(H_{i}\) , such that the map \((g,s)\mapsto gs\) induces by passage to the quotient an isomorphism from \(G\times^{H_{i}}S_{i}\) to \(U_{i}\) (show that the open subsets of X/G whose inverse image in X admits such a covering form a quasi-full subset of the set of open subsets of X/G, cf.~General Topology, Chap. IX, §4, Exerc. 27).

\begin{enumerate}
\def\labelenumi{\alph{enumi})}
\setcounter{enumi}{2}
\tightlist
\item
  With the assumptions in b), prove that there exists a representation of G on a finite dimensional real vector space V and an embedding of X in V of class \(C^{r}\) compatible with the operations of G (argue by induction on the set of subgroups of G, using b), Exerc. 2 and Exerc. 7).
\end{enumerate}

\begin{enumerate}
\def\labelenumi{\arabic{enumi})}
\setcounter{enumi}{8}
\tightlist
\item
  Let \(G\) be a Lie group operating properly on a manifold \(X\) . Make one of the following two assumptions:
\end{enumerate}

\begin{enumerate}
\def\labelenumi{(\roman{enumi})}
\tightlist
\item
  X/G is compact, and X is countable at infinity and of bounded dimension;\\
\item
  X is a finite dimensional vector space on which G operates linearly.
\end{enumerate}

Prove that the set of orbit types of elements of X is finite (treat the two cases simultaneously by induction on dim X, and use Prop. 6).

\[
\mathbf {G L} (3, \mathbf {R})
\]

\(\left(\begin{array}{ccc}1 & a & b \\ 0 & 1 & 0 \\ 0 & 0 & 1\end{array}\right), a, b \in \mathbf{R}.\) Show that, for the identity representation of \(G\) on

\(\mathbf{R}^3\) , the number of orbit types is infinite.

\end{enumerate}
